\documentclass[10pt,oneside]{article}
\usepackage[top=2cm,bottom=2cm,left=2cm,right=2cm]{geometry}
\usepackage{amssymb}
\usepackage{amsmath}
\usepackage{calc,multido,ifthen,color,bm}
\usepackage{theorem,graphicx}
\usepackage[frenchb]{babel}
\usepackage[T1]{fontenc}
\usepackage[utf8]{inputenc}
\usepackage{fancyhdr}
\usepackage{stmaryrd}
\usepackage{setspace}
\usepackage{dsfont}
\usepackage{hyperref}
\pagestyle{fancy}

%% DATE DE MISE A JOUR A MODIFIER ICI UNIQUEMENT : 
\def\datedemiseajour{01/10/2025}


\fancyhead[L]{Banque \'epreuve orale de math\'ematiques session 2026,  CCINP, filière MP et MPI}
\fancyhead[C]{}
\fancyhead[R]{Mise à jour : \datedemiseajour}
\fancyfoot[L]{}
\fancyfoot[C]{ \href{http://creativecommons.org/licenses/by-nc-sa/3.0/deed.fr}{CC BY-NC-SA 3.0 FR}}
\fancyfoot[R]{Page \thepage}

 \author{
} 
\date{2014,  \href{http://creativecommons.org/licenses/by-nc-sa/3.0/deed.fr}{CC BY-NC-SA 3.0 FR}}
%Le titre général de l'article centré verticalement
\title{\vspace{\fill}\textbf{ CONCOURS COMMUN INP \\ FILIÈRE  MP - FILIÈRE MPI \\\vspace{3cm}
\Huge{\textbf{BANQUE\\ \'EPREUVE ORALE\\DE MATH\'EMATIQUES\\SESSION 2026\\\vspace{3cm} sans corrigés }}}  \vspace{\fill}}

\begin{document}
\maketitle

\thispagestyle{empty}
\begin{flushright}
\textbf{Dernière mise à jour : le \datedemiseajour }
\end{flushright}




\pagebreak
\begin{flushleft}


%\vspace*{\stretch{1}}
\section*{Introduction}
L’épreuve orale de mathématiques du concours commun INP,  filière MP et filière MPI, se déroule de
la manière suivante :
\begin{itemize}
\item
25mn de préparation sur table.
\item 
25mn de passage à l’oral.
\end{itemize}
\bigskip

Chaque sujet proposé  est  constitué de deux exercices :
\begin{itemize}
\item
un exercice sur 8 points issu de la banque publique accessible sur le
site \url{https://www.concours-commun-inp.fr/fr/index.html}\\
\item
un exercice sur 12 points.
\end{itemize}

\medskip
Les deux exercices proposés portent sur des domaines différents.\\
\bigskip

Ce document contient les \textbf{112 exercices de la banque pour la
  session 2024}:
\begin{itemize}
\item
58 exercices d'analyse ( exercice 1 à exercice 58).\\
\item
36 exercices d'algèbre (exercice 59 à exercice 94). \\
\item
18 exercices de probabilités (exercice 95 à exercice 112).
\end{itemize}

\bigskip
Dans l'optique d'aider les futurs candidats à se préparer au mieux aux oraux du CCINP, chaque exercice de la banque  est proposé,  avec un corrigé, dans la version de la banque avec corrigés sur le site du CCINP.\\
\bigskip
Il se peut que des mises à jour aient lieu en cours d'année scolaire.\\
Cela dit, il ne s'agira, si tel est le cas, que de mises à jour mineures: reformulation de certaines questions pour plus de clarté, relevé d'éventuelles erreurs, suppression éventuelle de questions ou d'exercices.\\
Nous vous conseillons donc de vérifier, en cours d'année, en vous
connectant sur le site : 
\centerline{\url{https://www.concours-commun-inp.fr/fr/index.html}}
 Si une nouvelle version a été mise en ligne, la date de la dernière
 mise à jour figurant en haut  de chaque page.\\
Si tel est le cas, les exercices concernés et les modifications, seront signalés dans la version de la banque avec corrigés, sur le site du CCINP.

\bigskip
Remerciements à David DELAUNAY pour l'autorisation de libre utilisation du
fichier source de ses corrigés des exercices de l'ancienne banque, 
diffusés sur son site \url{http://mp.cpgedupuydelome.fr}

\vspace{1cm}
NB : la pr\'esente banque intègre des éléments issus des publications
suivantes :
\medskip
\begin{itemize}
\item[$\bullet$]
A. Antibi, L. d'Estampes et interrogateurs, Banque d'exercices de
mathématiques pour le programme 2003-2014 des oraux CCP-MP, {\it
  Éd. Ress. Pédag. Ouv. INPT}, {\bf 0701} (2013) 120 exercices.

\centerline{\url{http://pedagotech.inp-toulouse.fr/130701}}


\item[$\bullet$]
D. Delaunay, Prépas Dupuy de Lôme, cours et exercices corrigés MPSI -
MP, 2014.
\centerline{\url{http://mp.cpgedupuydelome.fr}}



\end{itemize}



\vspace{2cm}
\begin{flushright}
L'équipe des examinateurs de l'oral de mathématiques du CCINP, filière MP et filière MPI.
\end{flushright}

%\vspace*{\stretch{1}}
\vspace{.4cm}

\hfill
\begin{minipage}{0.45\textwidth}
\textbf{Contact}: {Valérie BELLECAVE, coordonnatrice des oraux
  de mathématiques du CCINP, filière MP et filière MPI.}
\\
\null\hfill \href{mailto:vbellecave@gmail.com}{vbellecave@gmail.com}
\end{minipage}


\end{flushleft}
\pagebreak
\begin{flushleft}

\part*{BANQUE ANALYSE}

\section*{EXERCICE 1\:analyse  }

On note $E$ l'espace vectoriel des applications continues sur $\left[ 0,1\right]$ à valeurs dans $\mathbb{R}$.\\
On pose : $\forall f\in E$,  $||f||_{\infty}=\underset{t\in \left[ 0,1\right] }{\sup}|f(t)|$ et $||f||_{1}=\displaystyle\int_{0}^{1}|f(t)|dt$.
\begin{enumerate}
\item
Les normes $||\:||_{\infty}$ et $||\:||_{1}$ sont-elles équivalentes? Justifier.
\item
Dans cette question, on munit $E$ de la norme $||\:||_{\infty}$.
\begin{enumerate}
\item
Soit $u:\left\lbrace \begin{array}{lll}
E&\longrightarrow& \mathbb{R}\\
f&\longmapsto& f(0)
\end{array}
\right. $\\
Prouver que $u$ est une application continue sur $E$.\:\:\:\:
\item
On pose $F=\left\lbrace  f\in E\:/\: f(0)=0\right\rbrace$.\\
Prouver que $F$ est une partie fermée de $E$  pour la norme $||\:||_{\infty}$.

\end{enumerate}
\item
Dans cette question, on munit $E$ de la norme $||\:||_{1}$.


Soit $c:\left\lbrace \begin{array}{lll}
\left[0,1 \right] &\longrightarrow& \mathbb{R}\\
x&\longmapsto& 1
\end{array}
\right. $\\
\medskip
On pose : $\forall n\in \mathbb{N}^{*}$, $f_n(x)=\left\lbrace \begin{array}{lll}
nx& \text{si}& 0\leqslant x\leqslant \dfrac{1}{n}\\
1&\text{si}& \dfrac{1}{n}< x\leqslant 1
\end{array}\right. $
\begin{enumerate}
\item
Soit $n\in \mathbb{N}^{*}$. Calculer $||f_n-c||_{1}$.
\item
On pose $F=\left\lbrace  f\in E\:/\: f(0)=0\right\rbrace$.\\
On note $\bar{F}$ l'adhérence de $F$.\\
Prouver que $c\in \bar{F}$\\
$F$ est-elle une partie fermée de $E$ pour la norme $||\:||_{1}$? 
\end{enumerate}
\end{enumerate}


\section*{EXERCICE 2 \:analyse }

On pose $f(x)=\dfrac{3x+7}{(x+1)^{2}}$~.
\begin{enumerate}
\item Décomposer $f(x)$ en éléments simples.

\item En déduire que $f$ est développable en série entière sur un intervalle du type $\left]-r,r \right[$ (où $r>0$).\\
Préciser ce développement en série entière et déterminer, en le justifiant, le domaine de validité $D$ de ce développement en série entière.\:\:\:\:\\

\item 
\begin{enumerate}
\item
Soit $\sum a_nx^n$ une série entière de rayon $R>0$.\\
On pose, pour tout $x\in \left]-R,R \right[ $, $g(x)=\displaystyle\sum_{n=0}^{+\infty}a_nx^n$.\\
Exprimer, pour tout entier $p$,  en le prouvant,  $a_p$ en fonction de $g^{(p)}(0)$.\\
\item
En déduire le développement limité de $f$ à l'ordre 3 au voisinage de 0.
\end{enumerate}

\end{enumerate}


\section*{EXERCICE 3\:analyse  }

\begin{enumerate}
\item
On pose $g(x)=\mathrm{e}^{2x}$ et $h(x)=\dfrac{1}{1+x}$.\\
Calculer, pour tout entier naturel $k$, la dérivée d'ordre $k$ des fonctions $g$ et $h$ sur leurs ensembles de définitions respectifs.
\item 
On pose $f(x)=\dfrac{\mathrm{e}^{2x}}{1+x}$.\\
En utilisant la formule de Leibniz concernant la dérivée $n^{\text{ième}}$ d'un produit de fonctions, déterminer, pour tout entier naturel $n$ et pour tout $x\in\mathbb{R}\backslash\left\lbrace -1\right\rbrace $, la valeur de $f^{(n)}(x)$.
\item 
Démontrer, dans le cas général, la formule de Leibniz, utilisée dans la question précédente.
\end{enumerate}

\section*{EXERCICE 4 \:analyse  }

\begin{enumerate}
\item \'Enoncer le théorème des accroissements finis.
\item Soit $f:\left[ a,b\right]\longrightarrow \mathbb{R} $ et 
soit $x_0\in \left]a,b \right[ $.\\
On suppose que $f$ est continue sur $[a,b]$ et  que $f$ est dérivable sur $]a,x_0[$ et sur $]x_0,b[$.\\
Démontrer que,  si $f'$ admet une limite finie en $x_0$,  alors 
 $f$ est dérivable en $x_0$ et $f'(x_0)=\lim\limits_{x\rightarrow x_0}f'(x)$.
\item Prouver que l'implication :\: ( $f$ est dérivable en $x_0$) $\Longrightarrow$ ($f'$ admet une limite finie en $x_0$) est fausse.

\textbf{Indication}: on pourra considérer la fonction $g$ définie par: $g(x)=x^2\sin\dfrac{1}{x}$ si $x\not=0$ et $g(0)=0$.
\end{enumerate}


\section*{EXERCICE 5 \:analyse }

\begin{enumerate}

\item
On considère la s\'{e}rie de terme g\'{e}n\'{e}ral $u_{n}=\dfrac{1}{n\left( \ln n\right) ^{\alpha }}$ o\`{u} $n\geqslant 2$ et $\alpha \in \mathbb{R}$.\medskip


\begin{enumerate}
\item \textbf{Cas $\bm{\alpha \leqslant 0}$}\medskip

En utilisant une minoration tr\`{e}s simple de $u_{n}$, démontrer que la série diverge.

\item \textbf{Cas $\bm{\alpha >0}$}\medskip 

\'Etudier la nature de la série.\\ 
\textbf{Indication}: on pourra utiliser la fonction $f$ d\'{e}finie par $f(x) =\dfrac{1}{x(\ln x) ^{\alpha }}$.
\end{enumerate}
\item
Déterminer la nature de la série $\displaystyle\sum\limits_{n\geqslant 2}^{}\dfrac{\left( \mathrm{e}-\left( 1+\dfrac{1}{n}\right)^n \right)\mathrm{e}^{\frac{1}{n}}}{\left( \ln(n^2+n)\right) ^2}$.
\end{enumerate}

\section*{EXERCICE 6 \:analyse }

Soit $\left( u_{n}\right) _{n\in \mathbb{N}}$ une suite de r\'{e}els strictement positifs et $l$ un r\'{e}el positif strictement inf\'{e}rieur \`{a} 1.

\begin{enumerate}
\item Démontrer que si $\underset{n\rightarrow +\infty }{\lim }\dfrac{u_{n+1}}{u_{n}}=l$, alors la s\'{e}rie $\displaystyle\sum u_{n}$ converge.
 
\textbf{Indication}: écrire, judicieusement, la définition de $\underset{n\rightarrow +\infty }{\lim }\dfrac{u_{n+1}}{u_{n}}=l$,  puis majorer, pour $n$ assez grand, $u_{n}$ par le terme g\'{e}n\'{e}ral d'une suite g\'{e}om\'{e}trique.
\item Quelle est la nature de la série $\displaystyle\sum_{n\geqslant1}{} \dfrac{n!}{n^{n}}$?
\end{enumerate}


\section*{EXERCICE 7 \:analyse }

\begin{enumerate}
\item
Soit  $\left( u_{n}\right)_{n\in \mathbb{N}}$ et $\left( v_{n}\right) _{n\in \mathbb{N}}$  deux suites réelles telles que $(v_n)_{n\in \mathbb{N}}$ est non nulle à partir d'un certain rang .

\begin{enumerate}
\item

Prouver que si  $u_{n}\underset{+\infty}\thicksim v_{n}$ alors $u_{n}$ et $v_{n}$ sont de m\^{e}me signe \`{a} partir d'un certain rang
\item
Dans cette question, on suppose que $(v_n)$ est positive.\\

Prouver que :
$u_{n}\underset{+\infty}{\thicksim} v_{n} \  \Longrightarrow  \ \displaystyle\sum u_{n}\ \text{et }\displaystyle\sum v_{n}\ \text{sont de m\^{e}me nature}.$

\end{enumerate}
\item 
Étudier la convergence de la s\'{e}rie
 $\displaystyle\sum\limits_{n\geqslant 2}^{} \dfrac{\left((-1)^n+\mathrm{i}\right) \sin \left( \dfrac{1}{n}\right) \ln n }{\left(\sqrt{n+3}-1\right)}$.


\textbf{Remarque 1} : $\mathrm{i}$ désigne le nombre complexe de carr\'{e} \'{e}gal \`{a} $-1$.\\

\end{enumerate}


\section*{EXERCICE 8 \:analyse }


\begin{enumerate}
\item 
Soit $\left( u_{n}\right) _{n\in \mathbb{N}}\ $une suite d\'{e}croissante positive de limite nulle.
\begin{enumerate}
\item Démontrer que la s\'{e}rie $\displaystyle\sum \left( -1\right) ^{k}u_{k}$ est convergente.

\textbf{Indication}: on pourra consid\'{e}rer $\left( S_{2n}\right) _{n\in \mathbb{N}}$ et $\left( S_{2n+1}\right) _{n\in \mathbb{N}}$ avec $S_{n}=\displaystyle\sum_{k=0}^{n}\left( -1\right) ^{k}u_{k}$.
\item Donner une majoration de la valeur absolue du reste de la série $\displaystyle\sum \left( -1\right) ^{k}u_{k}$.
\end{enumerate}
\item 
On pose : $\forall \:n\in\mathbb{N}^*$, $\forall\:x\in\mathbb{R}$, $f_n(x)=\dfrac{\left(-1\right)^{n}e^{-nx}}{n}$.
\begin{enumerate}

\item \'Etudier la convergence simple sur $\mathbb{R}$  de la série de fonctions $\displaystyle\sum_{n\geqslant1}^{}f_n$.
\item
\'Etudier la convergence uniforme sur $\left[ 0,+\infty\right[ $  de la série de fonctions $\displaystyle\sum_{n\geqslant1}^{}f_n$. 
\end{enumerate}
\end{enumerate}

\section*{EXERCICE 9 \:analyse }

\begin{enumerate}
\item 
 Soit $X$ un ensemble, $\left( g_{n}\right) $ une suite de fonctions de $X$ dans $\mathbb{C}$ et $g$ une fonction de $X$ dans $\mathbb{C}$.\\
Donner la définition de la convergence uniforme sur $X$ de la suite de fonctions $\left(g_n\right)$ vers la fonction $g$.\\

  \item   
On pose $f_{n}(x) =\dfrac{n+2}{n+1}\mathrm{e}^{-n x^{2}}\cos \left( \sqrt{n}x\right) $.
\begin{enumerate}
\item \'Etudier la convergence simple de la suite de fonctions $\left(f_{n}\right) $.
\item     La suite de fonctions  $\left(f_{n}\right)$ converge-t-elle uniformément sur $\left[ 0,+\infty\right[$ ?

\item Soit $a>0$. La suite de fonctions $\left(f_{n}\right) $ converge-t-elle uniform\'{e}ment sur  $[a,+\infty[$ ?	

\item La suite de fonctions  $\left(f_{n}\right)$ converge-t-elle uniform\'{e}ment sur $]0,+\infty[$? 
\end{enumerate}
\end{enumerate}
\section*{EXERCICE 10 \:analyse }

On pose $f_{n}\left( x\right) =\left( x^{2}+1\right) \dfrac{ne^{x}+xe^{-x}}{n+x}$.

\begin{enumerate}
\item Démontrer que la suite de fonctions $\left( f_{n}\right) $ converge uniform\'{e}ment sur $[0,1]$.
\item Calculer $\underset{n\rightarrow +\infty }{\lim }\displaystyle\int\limits_{0}^{1}\left( x^{2}+1\right) \dfrac{ne^{x}+xe^{-x}}{n+x}\text{d}x$.
\end{enumerate}



\section*{EXERCICE 11 \:analyse }

\begin{enumerate}
\item Soit $X$ une partie de $\mathbb{R}$, $\left( f_{n}\right) $ une suite de fonctions de $X$ dans $\mathbb{R}$ convergeant simplement vers une fonction $f$. \\
On suppose qu'il existe une suite $\left( x_{n}\right)_{n\in \mathbb{N}}$ d'\'{e}l\'{e}ments de $X$ telle que la suite $\left( f_{n}(x_{n})-f\left( x_{n}\right) \right) _{n\in \mathbb{N}}$ ne tende pas vers $0$. \bigskip

D\'{e}montrer que la suite de fonctions $\left( f_{n}\right) $ ne converge pas uniform\'{e}ment vers $f$ sur $X$.

\item Pour tout $x\in\mathbb{R}$, on pose $f_{n}(x) =\dfrac{\sin \left( nx\right) }{1+n^{2}x^{2}}$. 
	\begin{enumerate}
	\item \'Etudier la convergence simple de la suite $\left( f_{n}\right)$.
	\item \'Etudier la convergence uniforme de la suite $\left( f_{n}\right)$ sur $[a,+\infty[$ (avec $a>0$),  puis sur $]0,+\infty[$.
	\end{enumerate}
\end{enumerate}




\section*{EXERCICE 12\:analyse }

\begin{enumerate}
\item Soit $(f_n)$ une suite de fonctions de $[a,b]$ dans $\mathbb{R}$.

On suppose que la suite de fonctions $(f_n)$ converge uniformément sur $[a,b]$ vers une fonction $f$, et que, pour tout $n\in\mathbb{N}$, $f_n$ est continue en $x_0$, avec $x_0\in[a,b]$.

Démontrer que $f$ est continue en $x_0$.
\item On pose : $\forall\:n\in\mathbb{N}^*$, $\forall\:x\in[0;1]$, $g_n(x)=x^n$.\\
 La suite de fonctions $(g_n)_{n\in\mathbb{N}^*}$ converge-t-elle uniformément sur $[0;1]$?
\end{enumerate}


\section*{EXERCICE 13\:analyse  }


 \begin{enumerate}
  \item
  Rappeler, oralement, la définition, par les suites de vecteurs, d'une partie compacte d'un espace vectoriel normé.
  \item
  Démontrer qu'une partie compacte d'un espace vectoriel normé est une partie fermée de cet espace.
  \item
  Démontrer qu'une partie compacte d'un espace vectoriel normé est une partie bornée de cet espace.\\
 \textbf{ Indication }: On pourra raisonner par l'absurde.
  \item
  On se place sur $E=\mathbb{R}[X]$ muni de la norme $||\:||_{1}$ définie pour tout polynôme $P=a_0+a_1X+....+a_nX^n$ de $E$  par :
  $||P||_{1}=\displaystyle\sum_{i=0}^{n}|a_{i}|$.
  \begin{enumerate}
  
\item
  Justifier que $S(0,1)=\left\lbrace  P\in\mathbb{R}[X]\:/\: ||P||_{1}=1\right\rbrace $ est une partie fermée et bornée de $E$.
  \item
  Calculer $||X^n-X^m||_{1}$ pour $m$ et $n$ entiers naturels distincts.\\
  $S(0,1)$ est-elle une partie compacte de $E$? Justifier.
    \end{enumerate}
  \end{enumerate}

\section*{EXERCICE 14 \:analyse }

\begin{enumerate}
\item
Soit $a$ et $b$ deux réels donnés avec $a<b$.\\
  Soit $\left( f_{n}\right) $\ une suite de fonctions continues sur $[a,b],$ \`{a} valeurs r\'{e}elles. \\
Démontrer que si la suite $\left( f_{n}\right)$ converge uniform\'{e}ment sur $\left[ a,b\right] $ vers $f$, alors la suite $\left( \displaystyle\int_{a}^{b}f_{n}\left( x\right)\text{d}x\right)_{n\in \mathbb{N}}$ converge vers $\displaystyle\int_{a}^{b}f\left(x\right) \text{d}x$.

\item 
Soit $(f_n)$ une suite de fonctions continues sur $\left[ a,b\right] $.\\
On suppose que que la série $\displaystyle\sum f_n$ converge uniformément sur $\left[ a,b\right] $.\\
Prouver , \textbf{en utilisant 1.,} que $\displaystyle\int_{a}^{b}\left( \sum_{n=0}^{+\infty}f_n(x)\right) dx=\displaystyle\sum _{n=0}^{+\infty} \int_{a}^{b}f_n(x)dx$
\item
D\'{e}montrer que
$\displaystyle\int_{0}^{\frac{1}{2}}\left( \displaystyle\sum_{n=0}^{+\infty}x^{n}\right) \text{d}x=\displaystyle\sum\limits_{n=1}^{+\infty }\dfrac{1}{n2^{n}}~.$

\end{enumerate}

\section*{EXERCICE 15\:analyse }

Soit $X$ une partie de $\mathbb{R} $ ou $\mathbb{C} $ .
\begin{enumerate}
\item
Soit $\displaystyle\sum f_n$ une série de fonctions définies sur $X$ à valeurs dans $\mathbb{R}$ ou $\mathbb{C}$.\\
Rappeler la définition de la convergence normale de $\displaystyle\sum f_n$ sur $X$, puis celle  de la convergence uniforme de $\displaystyle\sum f_n$ sur $X$.

\item Démontrer que toute s\'{e}rie de fonctions, à valeurs dans $\mathbb{R} $ ou $\mathbb{C} $, normalement convergente sur $X$ est uniform\'{e}ment convergente sur $X$.

\item La s\'{e}rie de fonctions $\displaystyle\sum \dfrac{n^{2}}{n!}z^{n}$ est-elle uniform\'{e}ment convergente sur le disque ferm\'{e} de centre $0$ et de rayon $R\in \mathbb{R}_{+}^{\ast}$?
\end{enumerate}


\section*{EXERCICE 16 \:analyse }

  On consid\`{e}re la s\'{e}rie de fonctions de terme g\'{e}n\'{e}ral $u_{n}$ d\'{e}finie par: 
\begin{equation*}
\forall n\in \mathbb{N}^{*},\ \forall x\in \lbrack 0,1], \ \ u_{n}\left(x\right) =\ln \left( 1+\dfrac{x}{n}\right) -\dfrac{x}{n}~.
\end{equation*}

On pose, lorsque la s\'{e}rie converge, $S(x)=\displaystyle\sum\limits_{n=1}^{+\infty }\left[ \ \ln \left( 1+\dfrac{x}{n}\right) -\dfrac{x}{n}\right]$.

\begin{enumerate}
  \item D\'{e}montrer que $S$ est définie sur $[0,1]$. 

\item On définit une suite $(u_n)_{n\geq 1}$ par $u_n=\ln(n+1)-\displaystyle\sum_{k=1}^n\dfrac{1}{k}$.\\
  % \begin{enumerate}\item
  En utilisant $S(1)$ démontrer que la suite $(u_n)_{n\geq 1}$ est convergente.\\%  \:\:\:\:\bareme{1 point}
  % \item
  En déduire un équivalent simple de $\displaystyle\sum_{k=1}^n\dfrac{1}{k}$ lorsque $n\to+\infty$.

  % \end{enumerate}

\item D\'{e}montrer que $S$ est de classe $\mathcal{C}^1$ sur $[0,1]$ et calculer $S'(1)$. 

\end{enumerate}

\section*{EXERCICE 17 \:analyse }

Soit $A\subset \mathbb{C}$ et $\left( f_{n}\right)$ une suite de fonctions de $A$ dans $\mathbb{C}$.

\begin{enumerate}
\item D\'{e}montrer l'implication:
	\begin{eqnarray*}
	& \left( \text{la s\'{e}rie de fonctions }\displaystyle\sum f_{n}\ \text{converge uniform\'{e}ment sur $A$}\right)& \\
	&\Downarrow &\\
	&\left( \text{la suite de fonctions\ }\left( f_{n}\right)\ \text{converge uniform\'{e}ment vers 0 sur $A$}\right)&
	\end{eqnarray*}
\item
On pose: $\forall\:n\in\mathbb{N}$, $\forall\:x\in\left[ 0;+\infty\right[ $, $f_n(x)=nx^2\mathrm{e}^{-x\sqrt{n}}$.\\
Prouver que $\displaystyle\sum f_n$ converge simplement sur $\left[ 0;+\infty\right[$.\\
$\displaystyle\sum f_n$ converge-t-elle uniformément sur $\left[ 0;+\infty\right[$? Justifier.
\end{enumerate}

\section*{EXERCICE 18 \:analyse }

On pose: $\forall\:n\in\mathbb{N}^*$, $\forall\:x\in\mathbb{R}$, $u_n (x) =\dfrac{ ( - 1)^n x^n } { n}$.\\ 
On considère la série de fonctions $\displaystyle\sum\limits_{n\geqslant1}^{}u_n$.
\begin{enumerate}
\item \'Etudier la convergence simple de cette série.


On note $D$ l'ensemble des $x$ où cette série converge et $S(x)$ la somme de cette série pour $x\in D$. 

\item \begin{enumerate}
	\item La fonction $S$ est-elle continue sur $D$? 

	\item \'Etudier la convergence normale, puis la convergence uniforme de cette série sur $D$. \\

\item
 \'Etudier la convergence uniforme de cette série sur $\left[ 0,1\right] $.
	\end{enumerate}
\end{enumerate}

\section*{EXERCICE 19 \:analyse }

\begin{enumerate}
\item
\begin{enumerate}
\item
Justifier, oralement,  à l'aide du théorème de dérivation d'une série de fonctions, que la somme d'une série entière de la variable réelle est dérivable sur son intervalle ouvert de convergence.\\
\textbf{Remarque} : On pourra utiliser, sans le démontrer,   que  la série $\displaystyle\sum a_nx^n$ et la série $\displaystyle\sum na_nx^n$ ont même rayon de convergence.

\item
En déduire le développement en série entière à l'origine,  de la fonction de la variable réelle : $x\longmapsto \dfrac{1}{(1-x)^2}$.  
\end{enumerate}
\item
\begin{enumerate}
\item
Donner  le développement en série entière à l'origine de la fonction de la variable complexe : $z\longmapsto \dfrac{1}{1-z}$. \:
\item
Rappeler les résultats sur le produit de Cauchy de deux séries entières. \:\:\:\:
\item
En déduire le développement en série entière à l'origine, de la fonction de la variable complexe: $z\longmapsto \dfrac{1}{(1-z)^2}$.


\end{enumerate}
\end{enumerate}

\section*{EXERCICE 20 \:analyse }

\begin{enumerate}
\item Donner la définition du rayon de convergence d'une série entière de la variable complexe.


\item Déterminer le rayon de convergence de chacune des s\'{e}ries enti\`{e}res suivantes:

\begin{enumerate}
\item $\displaystyle\sum \dfrac{\left(n!\right)^{2}}{\left(2n\right)!}z^{2n+1}$. 


\item  $\displaystyle\sum n^{\left( -1\right)^{n}} z^{n}$.

\item
$\displaystyle\sum \cos nz^{n}$.
\end{enumerate}
\end{enumerate}

\section*{EXERCICE 21 \:analyse  }

\begin{enumerate}


\item  Donner la définition du rayon de convergence d'une série entière de la variable complexe.
\item
Soit $\left( a_{n}\right) _{n\in \mathbb{N}}$ une suite born\'{e}e telle que la s\'{e}rie $\displaystyle\sum a_{n}$ diverge.\\ Quel est le rayon de convergence de la s\'{e}rie enti\`{e}re $\displaystyle\sum a_{n}z^{n}$? Justifier.
\item Quel est le rayon de convergence de la s\'{e}rie enti\`{e}re $\displaystyle\sum \limits_{n\geqslant 1}^{}\left( \sqrt{n}\right)^{(-1)^n} \ln \left(1+\dfrac{1}{\sqrt{n}} \right)  z^{n}$? 
\end{enumerate}


\section*{EXERCICE 22 \:analyse }


\begin{enumerate}
\item Que peut-on dire du rayon de convergence de la somme de deux s\'eries enti\`{e}res? Le d\'emontrer.\\


\item D\'{e}velopper en s\'{e}rie enti\`{e}re au voisinage de $0$, en pr\'ecisant le rayon de convergence, la fonction $f\ :\ x\longmapsto \ln \left( 1+x\right)+\ln \left( 1-2x\right)$~.
\bigskip

La s\'{e}rie obtenue converge-t-elle pour $x=\dfrac{1}{4}$? $x=\dfrac{1}{2}$? $x=-\dfrac{1}{2}$?

 En cas de convergence, la somme de cette série est-elle continue en ces points ?
\\

\end{enumerate}


\section*{EXERCICE 23 \:analyse }

Soit $\left( a_{n}\right) _{n\in \mathbb{N}}$ une suite complexe telle que la suite $\left( \dfrac{\left\vert a_{n+1}\right\vert }{\left\vert a_{n}\right\vert }\right) _{n\in \mathbb{N}}$ admet une limite.

\begin{enumerate}
\item Démontrer que les s\'{e}ries enti\`{e}res $\displaystyle\sum a_{n}x^{n}$ et $\displaystyle\sum (n+1)a_{n+1}x^{n}$ ont le m\^{e}me rayon de convergence.

On le note $R$.

\item Démontrer que la fonction $x\longmapsto \displaystyle\sum\limits_{n=0}^{+\infty }a_{n}x^{n}$ est de classe $\mathcal{C}^1$  sur l'intervalle $]-R,R[$.
\end{enumerate}

\section*{EXERCICE 24 \:analyse }

\begin{enumerate}
\item Déterminer le rayon de convergence de la série entière $\displaystyle\sum\dfrac{x^n}{(2n)!}$~.\bigskip

On pose $S(x)=\displaystyle\sum_{n=0}^{+\infty}\dfrac{x^n}{(2n)!}$~.
\item Rappeler, sans démonstration,  le développement en série entière en 0 de la fonction $x\mapsto \text{ch}(x)$  et préciser le rayon de convergence.
\item \begin{enumerate}
	\item Déterminer $S(x)$.
	\item On considère la fonction $f$ définie sur $\mathbb{R}$ par:
	\begin{equation*}
	f(0)=1,\ \ f(x)=\text{ch}\sqrt x\text{ si $x>0$},\ \ f(x)=\cos\sqrt{-x}\text{ si $x<0$}\ .
	\end{equation*}
	Démontrer que $f$ est de classe $C^{\infty}$ sur $\mathbb{R}$.
	\end{enumerate}
\end{enumerate}

\section*{EXERCICE 25 \:analyse }

\begin{enumerate}
\item Démontrer que, pour tout entier naturel $n$, la fonction $t\longmapsto \dfrac{1}{1+t^{2}+t^{n}e^{-t}}$ est int\'{e}grable sur $[0,+\infty[$.
\item Pour tout $n\in\mathbb{N}$, on pose $u_{n}=\displaystyle\int_{0}^{+\infty }\dfrac{\text{d}t}{1+t^{2}+t^{n}e^{-t}}$. Calculer $\underset{n\rightarrow +\infty }{\lim }u_{n}$.
\end{enumerate}

		

\section*{EXERCICE 26 \:analyse }

Pour tout entier $n\geqslant 1$, on pose $I_{n}=\displaystyle\int_{0}^{+\infty } \dfrac{1}{\left( 1+t^{2}\right) ^n} \text{d}t$.

\begin{enumerate}
\item Justifier que $I_{n}$ est bien d\'{e}finie.
\item 
\begin{enumerate}
\item
\'Etudier la monotonie de la suite $\left( I_{n}\right) _{n\in\mathbb{N}^*}$.
\item
 D\'{e}terminer la limite de la suite $\left( I_{n}\right) _{n\in\mathbb{N}^*}$.
\end{enumerate}
\item La s\'{e}rie $\displaystyle\sum\limits_{n\geqslant 1}^{}(-1)^n I_{n}$ est-elle convergente ?

\end{enumerate}


\section*{EXERCICE 27 \:analyse }

Pour tout $n\in \mathbb{N}^*$, on pose\ $f_{n}\left( x\right) =\dfrac{e^{-x}}{%
1+n^{2}x^{2}}\ \ $et $u_{n}=\displaystyle\int_{0}^{1}f_{n}\left( x\right) \mathrm{d}x$.

\begin{enumerate}
\item \'Etudier la convergence simple de la suite de fonctions $\left( f_{n}\right) $  sur $[0,1]$.
\item
Soit $a\in\left] 0,1 \right[$.  La suite de  fonctions $\left( f_{n}\right) $ converge-t-elle uniformément sur $\left[a,1 \right]$? 
\item
 La suite de  fonctions $\left( f_{n}\right) $ converge-t-elle uniformément   sur $[0,1]$?
\item Trouver la limite de la suite $\left( u_{n}\right) _{n\in \mathbb{N}^*}.$
\end{enumerate}


\section*{EXERCICE 28 \:analyse }

\emph{N.B. : les deux questions sont ind\'{e}pendantes.}

\begin{enumerate}
\item La fonction $x\longmapsto \dfrac{e^{-x}}{\sqrt{x^{2}-4}}$ est-elle int\'{e}grable sur $]2,+\infty[$?
\item Soit $a$ un réel strictement positif.\\ La fonction $x\longmapsto \dfrac{\ln x}{\sqrt{1+x^{2a}}}$ est-elle int\'{e}grable sur $]0,+\infty[$?

\end{enumerate}

\section*{EXERCICE 29\:analyse }

On pose :  $\forall\:x\in ]0,+\infty[$,  $\forall\:t\in \left] 0,+\infty\right[ $, 
$f(x,t)=\text{e}^{-t}t^{x-1}$ .

\begin{enumerate}
\item Démontrer que : $\forall \:x\in \left]0,+\infty \right[$,  la fonction $t\mapsto f(x,t)$ est intégrable sur $\left] 0,+\infty\right[ $. \bigskip

On pose alors: $\forall \:x\in]0,+\infty[$, $\Gamma(x)=\displaystyle\int_0^{+\infty}\text{e}^{-t}t^{x-1}\text{d}t$.
\item Pour tout $x\in]0,+\infty[$, exprimer $\Gamma(x+1)$ en fonction de $\Gamma(x)$.
\item Démontrer que $\Gamma$ est de classe $C^1$ sur $]0;+\infty[$  et exprimer $\Gamma~\!'(x)$ sous forme d'intégrale.
\end{enumerate}


\section*{EXERCICE 30 \:analyse }

\begin{enumerate}
\item \'Enoncer le th\'{e}or\`{e}me de d\'{e}rivation sous le signe int\'{e}grale.
\item D\'{e}montrer que la fonction $f:x\longmapsto \displaystyle\int_{0}^{+\infty }e^{-t^{2}}\cos \left( xt\right) \text{d}t$ est de classe $C^{1}$ sur $\mathbb{R}$.
\item 
\begin{enumerate}
\item Trouver une \'{e}quation diff\'{e}rentielle lin\'{e}aire $\left(E\right)$ d'ordre $1$ dont $f$ est solution.
\item
 Résoudre $\left(E\right)$ .
 \end{enumerate}
\end{enumerate}



\section*{EXERCICE 31 \:analyse }

\begin{enumerate}

\item
Déterminer une primitive de $x\longmapsto \cos^4x$.
\item
R\'{e}soudre sur $\mathbb{R}$ l'\'{e}quation diff\'{e}rentielle: $y''+y=\cos^3 x$ en utilisant la m\'{e}thode de variation des constantes.
\end{enumerate}
\medskip


\section*{EXERCICE 32 \:analyse }


Soit l'équation différentielle: $x(x-1)y''+3xy'+y=0$.
\begin{enumerate}
\item Trouver les solutions de cette équation différentielle développables en série entière sur un intervalle $\left]-r,r \right[$ de $\mathbb{R}$, avec $r>0$.\\ Déterminer la somme des séries entières obtenues.

\item 
Est-ce que toutes les solutions de $x(x-1)y''+3xy'+y=0$ sur $\left]0;1 \right[$ sont les restrictions d'une fonction développable en série entière sur $\left]-1,1 \right[$?\\
\end{enumerate}


\section*{EXERCICE 33 \:analyse }
On pose : $\forall\:(x,y)\in\mathbb{R}^2\backslash \left\lbrace (0,0)\right\rbrace$,   $f\left( x,y\right) =\dfrac{xy}{\sqrt{x^{2}+y^{2}}}$ et $f\left( 0,0\right) =0$.

\begin{enumerate}
\item Démontrer que $f$ est continue sur $\mathbb{R}^{2}$.

\item Démontrer que $f$ admet des d\'{e}riv\'{e}es partielles en tout point de $\mathbb{R}^{2}$.
\item $f$ est-elle de classe $C^{1}$ sur $\mathbb{R}^{2}$? Justifier.

\end{enumerate}

\section*{EXERCICE 34 \:analyse }
  Soit $A$ une partie non vide d'un $\mathbb{R}$-espace vectoriel normé $E$.
\begin{enumerate}
\item
Rappeler la définition d'un point adhérent à $A$, en termes de voisinages ou de boules.
\item 
Démontrer que: $x\in\bar{A}\Longleftrightarrow\exists(x_n)_{n\in\mathbb N}$ telle que, $\forall n\in\mathbb{N}, x_n\in A$ et $\displaystyle\lim_{n\to+\infty}x_n=x$.


\item Démontrer que, si $A$ est un sous-espace vectoriel de $E$, alors $\bar{A}$ est un sous-espace vectoriel de $E$.

\item
Soient $B$ une autre partie non vide de $E$.
Montrer que $\overline{A\times B} = \bar{A}\times\bar{B}$.

\end{enumerate}


\section*{EXERCICE 35 \:analyse }



$E$ et $F$ désignent deux espaces vectoriels normés.\\
On note $||\:||_{E}$ ( respectivement $||\:||_{F}$) la norme sur $E$ (respectivement sur $F$).

\begin{enumerate}
\item Soient $f$ une application de $E$ dans $F$ et $a$ un point de $E$.

On considère les propositions suivantes:

\begin{description}
\item[P1.] $f$ est continue en $a$.
\item[P2.] Pour toute suite $(x_n)_{n\in\mathbb{N}}$ d'éléments de $E$ telle que $\displaystyle\lim_{n\to+\infty}x_n=a$, alors $\displaystyle\lim_{n\to+\infty}f(x_n)=f(a)$. 
\end{description}
Prouver que les propositions P1 et P2 sont équivalentes.

\item 
Soit $A$ une partie dense dans $E$, et soient $f$ et $g$ deux applications continues de $E$ dans $F$.

Démontrer que si, pour tout $x\in A$, $f(x)=g(x)$, alors $f=g$.
\end{enumerate}


\section*{EXERCICE 36 \:analyse }



Soient $E$ et $F$ deux espaces vectoriels normés sur le corps $\mathbb R$.\\
On note $||\:||_{E}$ ( respectivement $||\:||_{F}$) la norme sur $E$ (respectivement sur $F$).

\begin{enumerate}
\item Démontrer que si $f$ est une application linéaire de $E$ dans $F$, alors les propriétés suivantes sont deux à deux équivalentes:

	\begin{description}
	\item[P1.] $f$ est continue sur $E$.
	\item[P2.] $f$ est continue en $0_E$.
	\item[P3.] $\exists k>0$ tel que : $\forall x\in E, \left\Vert f(x)\right\Vert_{F} \leqslant k\left\Vert x\right\Vert_E$.
	\end{description}

\item Soit $E$ l'espace vectoriel des applications  continues de $[0;1]$ dans $\mathbb{R}$ muni de la norme définie par: $\Vert f\Vert_{\infty}=\sup\limits_{x\in[0;1]}|f(x)|$~.
On considère l'application $\varphi$ de $E$ dans $\mathbb{R}$ définie par: $\varphi(f)=\displaystyle\int_0^1 f(t)\text{d}t$.\bigskip

Démontrer que $\varphi$ est linéaire et continue.
%\item Pour toute fonction $f$ linéaire et continue de $E$ dans $F$, on pose $|||f|||=\displaystyle\sup_{\left\Vert x\right\Vert \leqslant1}\dfrac{\left\Vert f(x)\right\Vert}{\left\Vert x\right\Vert }\cdot$

%Démontrer que $f\longmapsto|||f|||$ est une norme dans l'espace vectoriel ${\cal L}(E,F)$ des applications linéaires et continues de $E$ dans $F$. 
\end{enumerate}



\section*{EXERCICE 37 \:analyse }

On note $E$ l'espace vectoriel des applications continues de $[0;1]$ dans $\mathbb{R}$.\\
 On pose: $\forall \:f\in E$,
$N_{\infty}(f)=\sup\limits_{x\in[0;1]} |f(x)|$ et $N_1(f)=\displaystyle\int_0^1|f(t)| \text{d}t$.


\begin{enumerate}
\item \begin{enumerate}
\item Démontrer que $N_{\infty}$ et $N_1$ sont deux normes sur $E$.
\item Démontrer qu'il existe $k>0$ tel que, pour tout $f$ de $E$, $N_1(f)\leq k N_{\infty}(f)$.
\item Démontrer que tout ouvert pour la norme $N_1$ est un ouvert pour la norme $N_{\infty}$.
\end{enumerate}
\item Démontrer que les normes $N_1$ et $N_{\infty}$ ne sont pas équivalentes.
\end{enumerate}

\section*{EXERCICE 38 \:analyse }


\begin{enumerate}
\item

On se place sur  $E=\mathcal{C} \left( [0,1],\mathbb{R}\right) $, muni de la norme $||\:||_1$ définie par  : $\forall\:f\in E$, $||f||_1=\displaystyle\int_{0}^{1}|f(t)|dt$.\\
Soit $u:\begin{array}{lll}
E &\longrightarrow &E\\
f&\longmapsto & u(f)=g
\end{array}$
 avec $\forall\: x\:\in [0,1]$, $g(x)=\displaystyle\int_{0}^{x}f(t)dt$.\\
 \medskip

On admet que $u$ est un endomorphisme de $E$.\\

Prouver que $u$ est continue et calculer $|||u|||$.\\
\textbf{Indication} : considérer, pour tout entier $n$ non nul, la fonction $f_n$ définie par  $f_n(t)=n\mathrm{e}^{-nt}$.

\item
Soit $n\in\mathbb{N}^*$. Soit $(a_1,a_2,...,a_n)\in\mathbb{R}^n$ un $n$-uplet\textbf{ non nul}, \textbf{fixé}.\\
\medskip
  Soit $u:\begin{array}{ccc}
 \mathbb{R}^n&\longrightarrow & \mathbb{R}\\
 (x_1,x_2,...,x_n)&\longmapsto&  \displaystyle\sum_{i=1}^{n}a_ix_i
 \end{array}$.\\

 \begin{enumerate}
\item 
Justifier que $u$ est continue quel que soit le choix de la norme sur $\mathbb{R}^n$.

   \item
On munit  $\mathbb{R}^n$ de $||\:||_{2}$ où $\forall x=(x_1,x_2,...,x_n)\in \mathbb{R}^n$, $||x||_2=\sqrt{\displaystyle\sum_{k=1}^{n}x_k^2}$.\\
 Calculer $|||u|||$.\\


 
 \end{enumerate}
 \item
Déterminer un espace vectoriel $E$, une norme sur $E$  et un endomorphisme de $E$ non continu pour la norme choisie. Justifier.


\end{enumerate}
\textbf{Remarque} : Les questions 1., 2. et 3. sont indépendantes.


\section*{EXERCICE 39 \:analyse }

 On note $l^2$ l'ensemble des suites $x=(x_n)_{n\in\mathbb{N}}$ de nombres réels telles que la série $\displaystyle\sum x_n^2$ converge.
\begin{enumerate}
\item
\begin{enumerate}
\item 
 Démontrer que, pour $x=(x_n)_{n\in\mathbb{N}} \in l^2$ et $y=(y_n) _{n\in\mathbb{N}}\in l^2$, la série $\displaystyle\sum x_ny_n$ converge.\\
	
On pose alors  $(x|y)=\displaystyle\sum_{n=0}^{+\infty} x_ny_n$.\\

\item 

  Démontrer que $l^2$ est un sous-espace vectoriel de l'espace vectoriel des suites de nombres réels.\\ 

	\end{enumerate}
	\medskip

Dans la suite de l'exercice, on admet que $(\:|\:)$ est un produit scalaire dans $l^2$.\\
 On suppose que $l^2$ est muni de ce produit scalaire et de la norme euclidienne associée,  notée $||\:||$.\\
\item



Soit  $p\in\mathbb{N}$. Pour tout $x=(x_n)\in l^2$, on pose $\varphi(x)=x_p$.\\
 Démontrer que $\varphi$ est une application linéaire et continue de $l^2$ dans $\mathbb{R}$.
 \item

 On considère l'ensemble $F$ des suites réelles presque nulles c'est-à-dire l'ensemble des suites réelles dont tous les termes sont nuls sauf peut-être un nombre fini de termes.\\
 Déterminer $F^{\perp}$ (au sens de $(\:|\:)$).\\
 Comparer $F$ et $\left( F^{\perp}\right) ^{\perp}$.
\end{enumerate}


\section*{EXERCICE 40 \:analyse }



Soit $A$ une algèbre de dimension finie admettant $e$ pour élément unité et munie d'une norme notée ||\:||.\\
On suppose que : $\forall (u,v)\in A^2$, $||u.v||\leqslant||u||.||v||$.\\

\begin{enumerate}
\item Soit $u$ un élément de $A$ tel que $\Vert u\Vert<1$.
	\begin{enumerate}
	\item Démontrer que la série $\displaystyle\sum u^n$ est convergente.
	\item Démontrer que $(e-u)$ est inversible et que $(e-u)^{-1}=\displaystyle\sum_{n=0}^{+\infty}u^n$.
	\end{enumerate}
\item Démontrer que, pour tout $u\in A$, la série $\displaystyle\sum\dfrac{u^n}{n!}$ converge.


\end{enumerate}





\section*{EXERCICE 41 \:analyse }

Soit $f$ l'application de $\mathbb{R}^2$ dans $\mathbb{R}$ d\'efinie par $f:(x,y)\mapsto 4x^2+12xy-y^2$. \\
\medskip
Soit $C=\{(x,y)\in\mathbb{R}^2,\ x^2+y^2=13\}$.

\medskip

\begin{enumerate}
\item
Justifier que $f$ atteint un maximum et un minimum sur $C$.
\item
Soit $(u,v)\in C$ un point o\`u $f$ atteint un de ses extremums.
\begin{enumerate}
\item
Justifier avec un th\'eor\`eme du programme qu'il existe un r\'eel $\lambda$ tel que le syst\`eme $(S)$ suivant soit v\'erifi\'e:\\
$(S)\ :\ \left\{\begin{array}{lll}4u&+6v&=\lambda u\\6u&-v&=\lambda v\end{array}\right.$
\item
Montrer que $(\lambda-4)(\lambda+1)-36=0$.\\
 En d\'eduire les valeurs possibles de $\lambda$.
\end{enumerate}
\item
D\'eterminer les valeurs possibles de $(u,v)$, puis donner le maximum et le minimum de $f$ sur $C$.
\end{enumerate}




\section*{EXERCICE 42 \:analyse }


\addcontentsline{toc}{subsection}{Exercice 2}
On considère les deux équations différentielles suivantes:\\
$   2xy'-3y=0 $\:\:\:\:\:\:  $(H)$\\
$   2xy'-3y=\sqrt{x} $\:\:\:  $(E)$\\

\begin{enumerate}
\item Résoudre l'équation $(H)$ sur l'intervalle $\left]  0,+\infty\right[ $.
\item Résoudre l'équation $(E)$ sur l'intervalle $\left]  0,+\infty\right[   $.

\item L'équation $(E)$ admet-elle des solutions sur l'intervalle $\left[ 0,+\infty\right[ $?
\end{enumerate}


\section*{EXERCICE 43 \:analyse }

 

Soit $x_0 \in \mathbb{R}$.\\
On définit la suite $(u_n)$ par $u_0=x_0$ et, $\forall n\in\mathbb{N}\:,\: u_{n+1}=\mathrm{Arctan}(u_n)$.
\begin{enumerate}
\item 
\begin{enumerate}

\item
 Démontrer que la suite $(u_n)$ est monotone et déterminer, en fonction de la valeur de $x_0$, le sens de variation de $(u_n)$.
\item Montrer que $(u_n)$ converge et déterminer sa limite.
\end{enumerate}
\item Déterminer l'ensemble des fonctions $h$, continues sur $\mathbb{R}$, telles que : $\forall x \in \mathbb{R}$, $h(x)=h(\mathrm{Arctan}\: x)$.
\end{enumerate}

\medskip

\section*{EXERCICE 44 \:analyse }
 

Soit $E$ un espace vectoriel normé. Soient $A$ et $B$ deux parties non vides de $E$. 
\begin{enumerate}
\item 
\begin{enumerate} 
\item Rappeler la caractérisation de l'adhérence d'un ensemble à l'aide des suites. 
\item Montrer que : $A \subset B \Longrightarrow \overline{A} \subset \overline{B}$.
\end{enumerate}
\item Montrer que : $\overline{A \cup B}=\overline{A} \cup \overline{B}$ \\
\textbf{ Remarque}: une réponse sans utiliser les suites est aussi acceptée.
\item \begin{enumerate}
\item Montrer que : $\overline{A \cap B} \subset \overline{A} \cap \overline{B}$.
\item Montrer, à l'aide d'un exemple, que l'autre inclusion n'est pas forcément vérifiée (on pourra prendre $E=\mathbb{R}$).
\end{enumerate}
\end{enumerate}


\section*{EXERCICE 45 \:analyse }


\textbf{ Les questions 1. et 2. sont indépendantes}.\\
Soit $E$ un $\mathbb{R}$-espace vectoriel normé.  On note $||\:||$ la norme sur $E$.\\

Soit $A$ une partie non vide de $E$.\\
 On note $\overline{A}$  l'adhérence de $A$.
\begin{enumerate}

\item 

\begin{enumerate}
\item
 Donner la caractérisation séquentielle de $\overline{A}$. 

\item Prouver que, si $A$ est convexe, alors $\overline{A}$ est convexe.
\end{enumerate}
\item 
On pose : $\forall x \in E,\:{d_A}(x) = \mathop {\inf }\limits_{a \in A} \left\| {x - a} \right\|$.
\begin{enumerate}

\item Soit $x \in{E}$. Prouver que $d_{A}(x) = 0 \Longrightarrow x \in \overline A $.
\item On suppose que $A$ est fermée et que :  $\forall (x,y) \in{E^2}$, $\forall t\in{[0,1]}$, $d_A(tx+(1-t)y)\leqslant td_A(x)+(1-t)d_A(y)$.\\
Prouver que $A$ est convexe.

\end{enumerate}
\end{enumerate}


\section*{EXERCICE 46 \:analyse }



On considère la série:$\displaystyle\sum\limits_{n\geqslant1}{}\cos \left( \pi \sqrt{n^2+n+1}\right) $. 
\begin{enumerate}
\item
Prouver que, au voisinage de $+\infty$, $\pi\sqrt{n^2+n+1}=n\pi+\dfrac{\pi}{2}+\alpha\dfrac{\pi}{n}+O\left(\dfrac{1}{n^2} \right)$ où $\alpha$ est un réel que l'on déterminera.
\item
En déduire que $\displaystyle\sum\limits_{n\geqslant 1}{}\cos \left( \pi \sqrt{n^2+n+1}\right) $ converge.
\item
 $\displaystyle\sum\limits_{n\geqslant 1}{}\cos \left( \pi \sqrt{n^2+n+1}\right) $ converge-t-elle absolument?

\end{enumerate}



\section*{EXERCICE 47 \:analyse }

Pour chacune des séries entières de la variable réelle suivantes, déterminer le rayon de convergence et calculer la somme de la série entière sur l'intervalle ouvert de convergence:\\
\begin{enumerate}
\item
$\displaystyle\sum_{n\geqslant 1}^{} \dfrac{3^nx^{2n}}{n}$.\:\:\:\:
\item
$\displaystyle\sum a_nx^n$ avec $\left\lbrace \begin{array}{l}
a_{2n}=4^n\\
a_{2n+1}=5^{n+1}
\end{array}\right. $\:\:\:\:

\end{enumerate}
\section*{EXERCICE 48 \:analyse }


$C^{0}\left( \left[ 0,1\right] ,\mathbb{R}\right) $ désigne l'espace vectoriel des fonctions continues sur  $
\left[ 0,1\right] $ à valeurs dans $\mathbb{R}$.\\

Soit $f\in C^{0}\left( \left[ 0,1\right] ,\mathbb{R}\right) $ telle que :  
$ \forall n\in \mathbb{N},\:
\displaystyle\int_{0}^{1}t^{n}f\left( t\right) \mathrm{d}t=0$.


\begin{enumerate}
\item \'Enoncer le th\'{e}or\`{e}me de Weierstrass d'approximation par des
fonctions polynomiales.
\item

Soit $\left( P_{n}\right) $ une suite de fonctions polynomiales
convergeant uniform\'{e}ment sur le segment $\left[ 0,1\right]$ vers $f$.
\begin{enumerate}
\item
 Montrer que la suite de fonctions $\left( P_{n}f\right) $
converge uniform\'{e}ment sur le segment $\left[ 0,1\right] $ vers $f^2$.

\item D\'{e}montrer que $\displaystyle\int_{0}^{1} f^2(t)
 \ \mathrm{d}t=\underset{n\rightarrow +\infty }{\lim }
\displaystyle\int_{0}^{1}P_{n}(t)f(t)\mathrm{d}t $.

\item  Calculer $\displaystyle\int
_{0}^{1}P_{n}\left( t\right) f\left( t\right) \mathrm{d}t$.
\end{enumerate}
\item  
En d\'{e}duire que $f$ est la fonction nulle sur le segment $\left[ 0,1
\right] .$

\end{enumerate}

\section*{EXERCICE 49 \:analyse }

Soit $\displaystyle\sum a_{n}$ une s\'{e}rie absolument convergente \`{a} termes
complexes. On pose $M=$ $\displaystyle\sum\limits_{n=0}^{+\infty }\left\vert
a_{n}\right\vert .$\\
 On pose : $ \forall n\in \mathbb{N},\: \forall t\in \left[
0,+\infty \right[,\: f_{n}\left( t\right) =\dfrac{a_{n}t^{n}}{n!}
\mathrm{e}^{-t}.$\\
\begin{enumerate}

\item  
\begin{enumerate}
\item
Justifier que la suite $(a_n)$ est bornée.
\item
Justifier que la série de fonctions   $\sum\limits_{}^{ }f_{n}$ converge simplement sur $\left[ 0,+\infty \right[$.\\
\medskip
On admettra, pour la suite de l'exercice, que  $f$: $t\mapsto \displaystyle\sum\limits_{n=0}^{+\infty }f_{n}(t)$ est continue sur $\left[ 0,+\infty \right[ .$
\end{enumerate}
\item
\begin{enumerate}
\item
Justifier que, pour tout $n\in{\mathbb{N}}$, la fonction $g_n:t\mapsto t^n\text{e}^{-t}$ est intégrable sur $\left[ 0,+\infty \right[$ et calculer $\displaystyle\int_{0}^{+\infty }g_n(t)  \mathrm{d}t$.\\
En déduire la convergence et la valeur de $\displaystyle\int_{0}^{+\infty }\left\vert
f_{n}\left( t\right) \right\vert \, \mathrm{d}t$.

\item Prouver que $\displaystyle\int_{0}^{+\infty }\left(
\displaystyle\sum\limits_{n=0}^{+\infty }\dfrac{a_{n}\,t^{n}}{n!}\,\mathrm{e}^{-t}\right) \mathrm{d}t=$ $
\displaystyle\sum\limits_{n=0}^{+\infty }a_{n}.$ 
\end{enumerate}
\end{enumerate}


\section*{EXERCICE 50 \:analyse }



On considère la fonction   $\;F : x \mapsto  \displaystyle\int_{0}^{+\infty} \dfrac{\text{e}^{-2\,t}}{x+t} \mathrm{d}t$.
\begin{enumerate}
\item
Prouver que  $F$ est définie et continue sur $\left]  0;+\infty\right[ $.
 \item 
Prouver que $x\longmapsto xF(x)$ admet une limite en $+\infty$ et déterminer la valeur de cette limite.

 \item
Déterminer un équivalent, au voisinage de $+\infty$,  de $F(x)$.

\end{enumerate}
 

\section*{EXERCICE 51 \:analyse }



\begin{enumerate}
\item Montrer que la série $\sum\dfrac{(2n)!}{(n!)^{2}2^{4n}(2n+1)}$ converge.

\medskip

On se propose de calculer la somme de cette série.

\medskip

\item Donner le développement en série entière en 0 de $t\longmapsto\dfrac{1}{\sqrt{1-t}}$ en précisant le rayon de convergence.\\
\textbf{Remarque}: dans l'expression du développement, on utilisera la notation factorielle.

\item En déduire le développement en série entière en 0 de $x\longmapsto\mathrm{Arcsin}\,x$ ainsi que son rayon de convergence.

\item En déduire la valeur de  $\displaystyle\sum\limits_{n=0}^{+\infty}\dfrac{(2n)!}{(n!)^{2}2^{4n}(2n+1)}$\,.

\end{enumerate}


\section*{EXERCICE 52 \:analyse }

Soit $\alpha\in\mathbb R$.\\ 
 On considère l'application définie sur $\mathbb{R}^2$ par  $f(x,y)=\begin{cases}
\dfrac{y^{4}}{x^{2}+y^{2}-xy} & \mathrm{si}\  (x,y)\not=(0,0)\\ 
\hfil\alpha & \mathrm{si}\  (x,y)=(0,0).
\end{cases} $
\begin{enumerate}
\item Prouver que : $\forall(x,y)\in\mathbb R^{2},\:x^{2}+y^{2}-xy\geqslant\dfrac{1}{2}(x^{2}+y^{2})$.

\item 

\begin{enumerate}
\item 
Justifier que le domaine de définition de $f$ est bien $\mathbb{R}^2$.

\item Déterminer $\alpha$ pour que $f$ soit continue sur $\mathbb R^{2}$.
\end{enumerate}
\item
Dans cette question, on suppose que $\alpha=0$.
\begin{enumerate}
\item
 Justifier l'existence de $\dfrac{\partial f}{\partial x}$ et $\dfrac{\partial f}{\partial y}$ sur $\mathbb R^{2}\setminus\big\{(0,0)\big\}$ et les calculer.

\item Justifier l'existence de $\dfrac{\partial f}{\partial x}(0,0)$ et $\dfrac{\partial f}{\partial y}(0,0)$ et  donner leur valeur.
\item $f$ est-elle de classe ${\cal C}^{1}$ sur $\mathbb R^{2}$? 

\end{enumerate}
\end{enumerate}

\section*{EXERCICE 53 \:analyse }

On considère, pour tout entier naturel $n$ non nul, la fonction $f_n$ définie sur $\mathbb{R}$ par $f_n(x)=\dfrac{x}{1+n^4x^4}$.
\begin{enumerate}
\item
\begin{enumerate}
\item
Prouver que $\displaystyle\sum\limits_{n\geqslant 1}^{}f_n$  converge simplement sur $\mathbb{R}$.\\
On pose alors : $\forall x\in\mathbb{R}$, $f(x)=\displaystyle\sum\limits_{n=1}^{+\infty}f_n(x)$.
\item
Soit $(a,b)\in\mathbb{R}^2$ avec $0<a<b$.\\
$\displaystyle\sum\limits_{n\geqslant 1}^{}f_n$ converge-t-elle normalement sur $\left[a,b \right]$? sur $\left[ a,+\infty\right[$?  
\item
$\displaystyle\sum\limits_{n\geqslant 1}^{}f_n$ converge-t-elle normalement sur $\left[ 0,+\infty\right[ $?
\end{enumerate}
\item
Prouver que $f$ est continue sur $\mathbb{R}^*$.
\item
Déterminer $\lim\limits_{x\to+\infty}^{}f(x)$.
\end{enumerate}
\medskip

\section*{EXERCICE 54 \:analyse }


Soit $E$ l'ensemble des suites à valeurs réelles qui convergent vers $0$.\\

\begin{enumerate}
\item
Prouver que $E$ est un sous-espace vectoriel de l'espace vectoriel des suites à valeurs réelles.
\item
On pose : $\forall\: u=(u_n)_{n\in\mathbb{N}}\in E$, $||u||=\sup\limits_{n\in\mathbb{N}}|u_n|$.
\begin{enumerate}
\item
Prouver que $||\:.\:||$ est une norme sur $E$.


\item
Prouver que :  $\forall\: u=(u_n)_{n\in\mathbb{N}}\in E$, $\displaystyle\sum\limits_{}^{}\dfrac{u_n}{2^{n+1}}$ converge.\\

\item
On pose : $\forall\: u=(u_n)_{n\in\mathbb{N}}\in E$, $f(u)=\displaystyle\sum\limits_{n=0}^{+\infty}\dfrac{u_n}{2^{n+1}}$.\\
Prouver que $f$ est continue sur $E$.

\end{enumerate}
\end{enumerate} 


\section*{EXERCICE 55 \:analyse }


Soit $a$ un nombre complexe.\\
On note $E$ l'ensemble des suites à valeurs complexes telles que :\\ $\forall\:n\in\mathbb{N}$, $u_{n+2}=2au_{n+1}+4(\mathrm{i}a-1)u_n$ avec $(u_0,u_1)\in  \mathbb{C} ^2$.
\begin{enumerate}
\item
\begin{enumerate}

\item 
Prouver que $E$ est un sous-espace vectoriel de l'ensemble des suites à valeurs complexes.
\item

Déterminer, en le justifiant, la dimension de $E$. 
\end{enumerate}
\item
Dans cette question, on considère la suite de $E$ définie par:
$u_0=1$ et $u_1=1$. \\
 Exprimer, pour tout entier naturel $n$, le nombre complexe  $u_n$ en fonction de $n$.\\
\medskip
\textbf{ Indication}: discuter suivant les valeurs de $a$. 
\end{enumerate}

\section*{EXERCICE 56 \:analyse }


  Soit $f$ la fonction définie sur $\mathbb{R}^2$ par : $\forall(x,y)\in\mathbb{R}^2$, $f(x,y)=2x^3+6xy-3y^2+2$.
  \begin{enumerate}
  \item
  $f$ admet-elle des extrema locaux sur $\mathbb{R}^2$?
  Si oui, les déterminer.
  \item
  $f$ admet-elle des extrema globaux sur $\mathbb{R}^2$? Justifier.
  \item
  On pose $K=\left[ 0,1\right]\times \left[0,1 \right]$. \\
  Justifier, oralement,  que $f$ admet un maximum global sur $K$ puis  le déterminer.
  \end{enumerate}
 


\section*{EXERCICE 57 \:analyse }


\begin{enumerate}
 \item Soit $f $ une fonction de $\mathbb{R}^{2}$ dans $\mathbb{R}$.
 \begin{enumerate}
 \item Donner, en utilisant des quantificateurs, la définition de la continuité de $f$ en $(0,0)$.
 \item Donner la définition de «$f$ différentiable en  $(0,0)$».
 \end{enumerate}
 \item 
On considère l'application définie sur $\mathbb{R}^2$ par 
 $f(x,y)=\left\lbrace
  \begin{array}{ll}
 xy\dfrac{x^2-y^2}{x^2+y^2}& \text{si}\: (x,y)\neq (0,0)\\
 0 &\text{si} \:(x,y)=(0,0)
 \end{array}
 \right. $

 \begin{enumerate}
 \item Montrer que $f$ est continue sur $\mathbb{R}^2$.
 \item Montrer que $f$ est de classe $\mathcal {C}^1$ sur $\mathbb{R}^2$.
 \end{enumerate}
 \end{enumerate}
   

 
\section*{EXERCICE 58 \:analyse }

 \begin{enumerate}
 \item
 Soit $E$ et $F$ deux $\mathbb{R}$-espaces vectoriels normés de dimension finie.\\
  Soit $a\in E$ et soit $f:E \longrightarrow F$ une application.\\

 Donner la définition de «$f$ différentiable en  $a$».
 \item
 Soit $n\in\mathbb{N}^*$. Soit $E$ un $\mathbb{R}$-espace vectoriel de dimension finie $n$.\\
 Soit $e=(e_1,e_2,\dots,e_n)$ une base de $E$.\\
 On  pose :  $\forall x\in E,\; \|x\|_{\infty}=\underset{1\leqslant i\leqslant n}{\max}|x_i|,\mbox{ o\`u }x=\displaystyle\sum_{i=1}^{n}x_ie_i$.\\
On pose :  $ \forall (x,y)\in E\times E,\; \|(x,y)\|=\max(\|x\|_\infty,\|y\|_\infty)$.

\medskip

 On admet que $\|.\|_{\infty}$  est une norme sur $E$ et que $\|.\|$ est une norme sur $E\times E$.\\
 \medskip
 Soit $B:E\times E\longrightarrow \mathbb{R}$ une forme bilinéaire sur $E$.\\
 \begin{enumerate}
 \item
  Prouver que $\exists\: C\in\mathbb{R}^+/\:\forall\:(x,y)\in E\times E,\:|B(x,y)|\leqslant C\|x\|_{\infty}\|y\|_{\infty}$.
  \item
   Montrer que $B$ est différentiable sur $E\times E$ et déterminer sa différentielle en tout $(u_0,v_0)\in E\times E$.
  
  \end{enumerate}
\end{enumerate}

 \pagebreak
 \part*{BANQUE ALGÈBRE}
 \section*{EXERCICE 59 \: algèbre }
Soit $n$ un entier naturel tel que $n\geqslant 2$.\\
Soit $E$ l'espace vectoriel des polyn\^{o}mes \`{a} coefficients dans $\mathbb{K}$ ($\mathbb{K}=\mathbb{R}$ ou $\mathbb{K}=\mathbb{C}$) de degr\'{e} inf\'{e}rieur ou \'{e}gal \`{a} $n$.\\
On pose : $\forall\:P\in E$,  $f\left( P\right)=P-P^{\prime }$.

\begin{enumerate}
\item D\'{e}montrer que $f$ est bijectif de $E$ dans $E$ de deux mani\`{e}res:
	\begin{enumerate}
	\item sans utiliser de matrice de $f$, 
	\item en utilisant une matrice de $f$.  
	\end{enumerate}
	
\item Soit $Q\in E.$ Trouver $P$ tel que $f\left( P\right) =Q$~.

\textbf{Indication }: si $P\in E$, quel est le polyn\^{o}me $P^{\left(n+1\right) }$? 
\item
$f$ est-il diagonalisable?

\end{enumerate}

\section*{EXERCICE 60 \: algèbre }

Soit la matrice $A=\left( 
\begin{array}{cc}
1 & 2 \\ 
2 & 4%
\end{array}%
\right) \ \ $et $f$ l'endomorphisme de $\mathcal{M}_{2}\left( \mathbb{R}%
\right) $ d\'{e}fini par :\ \ \ $f\left( M\right) =AM.$

\begin{enumerate}

%\marginpar{3pt}
\item D\'{e}terminer une base de  $\mathrm{Ker}f$. 

%\marginpar{2pt}
\item $f$ est-il surjectif ? 

%\marginpar{3pt}
\item Déterminer une base de $ \text{Im}f.$\smallskip
\item
A-t-on $\mathcal{M}_{2}\left( \mathbb{R}%
\right)=\mathrm{Ker}f\oplus \mathrm{Im}f$?
\end{enumerate}

\section*{EXERCICE 61 \: algèbre }

On note $\mathcal{M}_{n}\left( \mathbb{C}\right)$ l'espace vectoriel des matrices carr\'{e}es d'ordre $n$ \`{a} coefficients complexes.

Pour $A=\left( a_{i,j}\right) _{\substack{ 1\leqslant i\leqslant n  \\ 1\leqslant j\leqslant n}}\in \mathcal{M}_{n}\left( \mathbb{C}\right)$, on pose : 
$\left\Vert A\right\Vert =\underset{_{\substack{ 1\leqslant i\leqslant n  \\ 1\leqslant j\leqslant n}}}{\max }\left\vert a_{i,j}\right\vert$.

\begin{enumerate}


\item
Prouver que $||\:||$ est une norme sur $\mathcal{M}_{n}\left( \mathbb{C}\right)$.
\item Démontrer que: $\forall\:(A,B)\in \left( \mathcal{M}_{n}\left( \mathbb{C}\right)\right) ^2$,  $\left\Vert AB\right\Vert \leqslant n\left\Vert A\right\Vert \left\Vert B\right\Vert $.\\
Puis, démontrer que,  pour tout entier $p\geqslant 1$, $\left\Vert A^{p}\right\Vert \leqslant n^{p-1}\left\Vert A\right\Vert ^{p}$.

\item Démontrer que, pour toute matrice $A\in \mathcal{M}_{n}\left( \mathbb{C}\right)$, la s\'{e}rie $\displaystyle\displaystyle\sum \dfrac{A^{p}}{p!}$ est absolument convergente.

Est-elle convergente?
\end{enumerate}
		


\section*{EXERCICE 62 \: algèbre }

Soit $E\ $un espace vectoriel sur $\mathbb{R}$ ou $\mathbb{C}$.\\
Soit $f\in\mathcal{L}(E)$ tel que $f^2-f-2\mathrm{Id}=0$.\\

\begin{enumerate}
\item 
 Prouver que $f$ est bijectif et exprimer $f^{-1}$ en fonction de $f$.\\
\item
Prouver que $E=\mathrm{Ker} (f+\mathrm{Id})\oplus\mathrm{Ker}(f-2\mathrm{Id})$:
\begin{enumerate}
\item
en utilisant le lemme des noyaux.
\item
sans utiliser le lemme des noyaux.
\end{enumerate}
\item
Dans cette question, on suppose que $E$ est de dimension finie.\\
Prouver que $\mathrm{Im}(f+\mathrm{Id})=\mathrm{Ker}(f-2\mathrm{Id})$.
\end{enumerate}
\section*{EXERCICE 63 \: algèbre }


Soit $E$ un espace euclidien muni d'un produit scalaire noté $(\:|\:)$.\\
On pose $\forall x\in E$, $||x||=\sqrt{(x|x)}$.\\
Pour tout endomorphisme $u$ de $E$, on note $u^*$ l'adjoint de $u$.
\begin{enumerate}
\item

Un endomorphisme $u$ de $E$ vérifiant $\forall x\in E$, $(u(x)|x)=0$ est-il nécessairement l'endomorphisme nul?
\item
Soit $u\in \mathcal{L}(E)$.\\
Prouver que les trois assertions suivantes sont équivalentes:\\

i.\: $u\circ u^*=u^*\circ u$.\\
ii.\: $\forall (x,y)\in E^2$, $(u(x)|u(y))=(u^*(x)|u^*(y))$.\\
iii.\: $\forall x\in E$, $||u(x)||=||u^*(x)||$.
\end{enumerate}








\section*{EXERCICE 64 \: algèbre }


Soit $f$ un endomorphisme d'un espace vectoriel $E$ de dimension finie $n$.
\begin{enumerate}
\item Démontrer que: $E=\text{Im} f \oplus \mathrm{Ker} f \Longrightarrow \text{Im} f = \text{Im} f^2$.\:\:\:\:

\item \begin{enumerate}
	\item Démontrer que: $\text{Im} f = \text{Im} f^2 \Longleftrightarrow \mathrm{Ker} f = \mathrm{Ker} f^2$.\:\:\:\:
	\item Démontrer que: $\text{Im} f = \text{Im} f^2 \Longrightarrow E=\text{Im} f \oplus \mathrm{Ker} f$.\:\:\:\:
	\end{enumerate}
\end{enumerate}
\section*{EXERCICE 65 \: algèbre }

Soit $u$ un endomorphisme d'un espace vectoriel $E$ sur le corps $\mathbb{K}$ ($=\mathbb{R}$ ou $\mathbb{C}$). On note $\mathbb{K}[X]$ l'ensemble des polynômes à coefficients dans $\mathbb{K}$.
\begin{enumerate}
\item Démontrer que : 
\begin{equation*}
\forall (P,Q)\in \mathbb{K}[X]\times \mathbb{K}[X],~ (PQ)(u)=P(u)\circ Q(u)~.
\end{equation*}
\item \begin{enumerate}
	\item Démontrer que : $\forall (P,Q)\in \mathbb{K}[X]\times \mathbb{K}[X],~ P(u)\circ Q(u)=Q(u)\circ P(u)$~.
	\item Démontrer que, pour tout $(P,Q)\in \mathbb{K}[X]\times \mathbb{K}[X]$:
	\begin{equation*}
	\text{($P$ polynôme annulateur de $u$)$\Longrightarrow$ ($PQ$ polynôme annulateur de $u$)}
	\end{equation*}
	\end{enumerate}
\item Soit $A=\begin{pmatrix}
-1 & -2 \\ 
1 & 2
\end{pmatrix} $. Écrire le polynôme caractéristique de $A$, puis en déduire que le polynôme $R=X^4+2X^3+X^2-4X$ est un polynôme annulateur de $A$.
\end{enumerate}

		

\section*{EXERCICE 66 \: algèbre}


\begin{enumerate}
\item
Soit $A\in S_n\left( \mathbb{R}\right)$.\\
Prouver que $A\in S_n^+\left( \mathbb{R}\right)\Longleftrightarrow \mathrm{sp}(A)\subset \left[0,+\infty \right[$. 
\item
Prouver que $\forall\:A\in S_n\left( \mathbb{R}\right)$, $A^2\in S_n^+\left( \mathbb{R}\right)$.


\item
Soit $A\in S_n^+\left( \mathbb{R}\right)$.\\
Prouver qu'il existe $B\in S_n^+\left( \mathbb{R}\right)$ telle que $A=B^2$.

\end{enumerate}




\section*{EXERCICE 67 \: algèbre}

Soit la matrice $M=\begin{pmatrix}
0 & a & c \\ 
b & 0 & c \\ 
b & -a & 0
\end{pmatrix}$ o\`{u} $a,b,c$ sont des r\'{e}els.


 $M$ est-elle diagonalisable dans $\mathcal{M}_{3}\left(\mathbb{R}\right)$?
 $M$ est-elle diagonalisable dans $\mathcal{M}_{3}\left( \mathbb{C}\right)$?



\section*{EXERCICE 68 \: algèbre }

Soit la matrice $A=\begin{pmatrix}
1 & -1 & 1 \\ 
-1 & 1 & -1 \\ 
1 & -1 & 1
\end{pmatrix}$~.
\begin{enumerate}
\item Démontrer que $A$ est diagonalisable de trois manières:
	\begin{enumerate}
	\item sans calcul, \:\:\:\:
	\item en calculant directement le déterminant $\text{det}(\lambda \mathrm{I}_3-A)$, où $\mathrm{I}_3$ est la matrice identité d'ordre 3, et en déterminant les sous-espaces propres,\:\:\:\:

	\item en calculant $A^2$.\:\:\:\:
	\end{enumerate}
\item
On note $f$ l'endomorphisme de $\mathbb{R}^3$ dont la matrice dans la base canonique de $\mathbb{R}^3$ est A.\\

 Trouver une base orthonormée dans laquelle la matrice de $f$ est diagonale. 

\end{enumerate}

\section*{EXERCICE 69 \: algèbre }

On considère la matrice $A=\begin{pmatrix}
0&a&1\\
a&0&1\\
a&1&0
\end{pmatrix}$ où $a$ est un réel.
\begin{enumerate}
\item
Déterminer le rang de $A$.
\item
Pour quelles valeurs de $a$, la matrice $A$ est-elle diagonalisable? 
\end{enumerate}


\section*{EXERCICE 70 \: algèbre  }

Soit $A=\begin{pmatrix}
0 & 0 & 1 \\ 
1 & 0 & 0 \\ 
0 & 1 & 0
\end{pmatrix} \in \mathcal{M}_{3}\left( \mathbb{C}\right)$~.

\begin{enumerate}
\item D\'{e}terminer les valeurs propres et les vecteurs propres de $A$. $A$ est-elle diagonalisable?
\item Soit $\left( a,b,c\right) \in \mathbb{C}^{3}$ et $B=a\mathrm{I}_{3}+bA+cA^{2}$, où $\mathrm{I}_3$ désigne la matrice identité d'ordre 3.\\ D\'{e}duire de la question \textbf{1.} les \'{e}l\'{e}ments propres de $B$.
\end{enumerate}

\section*{EXERCICE 71 \: algèbre}
Soit  $P$ le plan d'équation $x+y+z=0$ et  $D$ la droite d'équation $x=\dfrac{y}{2}=\dfrac{z}{3}$.
\begin{enumerate}
\item
Vérifier que $\mathbb{R}^3=P\oplus D$.\:\:\:\:
\item
Soit $p$ la projection vectorielle de  $\mathbb{R}^3$ sur  $P$  parallèlement à  $D$.\\
Soit $u=(x,y,z)\in\mathbb{R}^3$.\\
Déterminer $p(u)$ et donner la matrice de $p$ dans la base canonique de $\mathbb{R}^3$.\:\:\:\:
\item
Déterminer une base de $\mathbb{R}^3$ dans laquelle la matrice de $p$ est diagonale.\:\:\:\:

\end{enumerate}


\section*{EXERCICE 72\:algèbre }
Soit $n$ un entier naturel non nul.\\
Soit $f$ un endomorphisme d'un espace vectoriel $E$ de dimension $n$, et soit $e=\left( e_1,\ldots,e_n\right) $ une base de $E$.

On suppose que  $f(e_1)=f(e_2)=\cdots=f(e_n)=v$, où $v$ est un vecteur donné de $E$.
\begin{enumerate}
\item 
Donner le rang de $f$.

\item 
$f$ est-il diagonalisable? (discuter en fonction du vecteur $v$)
\end{enumerate}
		



\section*{EXERCICE 73 \: algèbre}

On pose $A=\begin{pmatrix}
2 & 1 \\ 
4 & -1
\end{pmatrix}$.

\begin{enumerate}

\item  Déterminer les valeurs propres et les vecteurs propres de $A$.
\item Déterminer toutes les matrices qui commutent avec la matrice $\begin{pmatrix}
3 & 0 \\ 
0 & -2
\end{pmatrix}$.\\
En déduire que l'ensemble des matrices qui commutent avec $A$ est $\mathrm{Vect}\left( \mathrm{I}_2,A\right) $.
\end{enumerate}

  \section*{EXERCICE 74 \: algèbre }



\begin{enumerate}
\item On considère la matrice $A=\begin{pmatrix}
1 & 0 & 2 \\ 
0 & 1 & 0\\ 
2 & 0 & 1
\end{pmatrix} $.
	\begin{enumerate}
	\item Justifier sans calcul que $A$ est diagonalisable.
	\item Déterminer les valeurs propres de $A$ puis une base de vecteurs propres associés.
	
	\end{enumerate}
\item On considère le système différentiel $\left\{\begin{array}{l}
 x'=x+2z\\ 
 y'=y \\ 
 z'=2x+z  
\end{array}\right.$ , $x,y,z$ désignant trois fonctions de la variable $t$, dérivables sur $\mathbb{R}$.

En utilisant la question 1. et en le justifiant, résoudre ce système.
\end{enumerate}

\section*{EXERCICE 75 \: algèbre }

On considère la matrice $A=\begin{pmatrix}
-1 & -4 \\ 
1 & 3
\end{pmatrix}$~.

\begin{enumerate}
\item Démontrer que $A$ n'est pas diagonalisable.

\item On note $f$ l'endomorphisme de $\mathbb{R}^{2}$ canoniquement associ\'{e} \`{a} $A$. \\Trouver une base $\left( v_{1},v_{2}\right)$ de $\mathbb{R}^{2}$ dans laquelle la matrice de $f$ est de la forme $\begin{pmatrix}
a & b \\ 
0 & c
\end{pmatrix}$.\\
On donnera explicitement les valeurs de $a$, $b$ et $c$.

\item En déduire la résolution du syst\`{e}me différentiel $\left\{
\begin{array}{l}
x^{\prime }=-x-4y \\ 
y^{\prime }=x+3y
\end{array}
\right.$~.
\end{enumerate}
		
	

\section*{EXERCICE 76 \: algèbre  }

Soit $E$ un $\mathbb{R}$-espace vectoriel muni  d'un produit scalaire noté $(\:|\:)$.\\
On pose $\forall\:x\in E$, $||x||=\sqrt{(x|x)}$.
\begin{enumerate}
\item 
\begin{enumerate}
\item 
\'Enoncer et d\'{e}montrer l'in\'{e}galit\'{e} de Cauchy-Schwarz.



\item
Dans quel cas a-t-on égalité? Le démontrer.
\end{enumerate}
\item

Soit $E=\left\lbrace f\in\mathcal{C}\left( \left[ a,b\right] ,\mathbb{R}\right) ,\:\forall\:x\in \left[ a,b\right]\:f(x)>0  \right\rbrace$. \\
Prouver que l'ensemble   $ \left\lbrace \displaystyle\int_{a}^{b}f(t)\mathrm{d}t\times \displaystyle\int_{a}^{b}\dfrac{1}{f(t)}\mathrm{d}t\:,\:f\in E\right\rbrace $ admet une borne inférieure $m$ et déterminer la valeur de $m$.
\end{enumerate}

\section*{EXERCICE 77 \: algèbre }


Soit $E$ un espace euclidien.

\begin{enumerate}


\item 
 Soit $A$ un sous-espace vectoriel de $E$.\\
D\'{e}montrer que $\left( A^{\perp }\right) ^{\perp }=A$.
\item
Soient  $F$ et $G$ deux sous-espaces vectoriels de $E$.
\begin{enumerate}

\item

D\'{e}montrer que \ $\left( F+G\right) ^{\perp }=F^{\perp }\cap G^{\perp }$.

\item D\'{e}montrer que\ $\left( F\cap G\right) ^{\perp }=F^{\perp}+G^{\perp }$.
\end{enumerate}
\end{enumerate}		



\section*{EXERCICE 78 \: algèbre  }

Soit $E$ un espace euclidien de dimension $n$ et $u$ un endomorphisme de $E$.\\ On note $\left( x|y\right) $ le produit scalaire de $x$ et de $y$ et $||.||$ la norme euclidienne associée.

\begin{enumerate}
%\item Démontrer que si :$\left( \forall (x,y)\in E^{2}\right)\ \left( u(x)|u(y)\right) =\left( x|y\right) $ alors $u$ est bijectif.
\item Soit $u$ un endomorphisme de $E$, tel que: $\forall x \in E, \vert\vert u(x)\vert\vert = \vert\vert x\vert\vert$. 
	\begin{enumerate}
	\item Démontrer que: $\forall (x,y)\in E^{2}~(u(x)|u(y)) = (x|y) $.
	\item Démontrer que $u$ est bijectif.
	\end{enumerate}
\item
On note ${\mathcal{O}}(E)$ l'ensemble des isométries vectorielles de $E$.\\
C'est -à-dire  ${\mathcal{O}}(E)=\left\lbrace u\in\mathcal{L}\left(E \right)\:/ \: \forall x \in E, \vert\vert u(x)\vert\vert = \vert\vert x\vert\vert \right\rbrace $.\\
\medskip
Démontrer que ${\mathcal{O}}(E)$ , muni de la loi $\circ $ , est un groupe.\\
\item
Soit $u\in\mathcal{L}(E)$. Soit $e=(e_1,e_2,...,e_n)$ une base orthonormée de $E$.\\
Prouver que : $u\in {\mathcal{O}}(E)\Longleftrightarrow$$ \left( u(e_1),u(e_2),...,u(e_n)\right) $ est une base orthonormée de $E$.
\end{enumerate}


\section*{EXERCICE 79 \: algèbre}

Soit $a$ et $b$ deux réels tels que $a$<$b$.\\
\begin{enumerate}
\item Soit $h$ une fonction continue et positive de $[a,b]$ dans $\mathbb{R}$.\bigskip

Démontrer que  $\displaystyle\int_{a}^{b}h(x)\text{d}x=0\Longrightarrow h=0$~.

\item Soit $E$ le $\mathbb{R}$-espace vectoriel des fonctions continues de $[a,b]$ dans $\mathbb{R}$.\\
  On pose : $\forall\:(f,g)\in E^2$, $\left( f|g\right) =\displaystyle\displaystyle\int_{a}^{b}f(x)g(x)\text{d}x$.\\
   Démontrer que l'on d\'{e}finit ainsi un produit scalaire sur $E$.

\item Majorer $\displaystyle\int_{0}^{1}\sqrt{x}e^{-x}\text{d}x$ en utilisant l'in\'{e}galit\'{e} de Cauchy-Schwarz.
\end{enumerate}


\section*{EXERCICE 80 \: algèbre}

Soit $E\ $l'espace vectoriel des applications continues et $2\pi $-p\'{e}riodiques de $\mathbb{R}$ dans $\mathbb{R}$.

\begin{enumerate}
\item Démontrer que $\left( f\ |\ g\right) =\dfrac{1}{2\pi }\displaystyle\int_{0}^{2\pi }f\left( t\right) g\left( t\right) \text{d}t$ d\'{e}finit un produit scalaire sur $E$.

\item Soit $F$ le sous-espace vectoriel engendr\'{e} par $f:x\mapsto \cos x$ et $g:x\mapsto \cos \left( 2x\right)$.\bigskip

D\'{e}terminer le projet\'{e} orthogonal sur $F$ de la fonction $u:x\mapsto \sin ^{2}x$.
\end{enumerate}


\section*{EXERCICE 81 \: algèbre}

On d\'{e}finit dans $\mathcal{M}_{2}\left( \mathbb{R}\right) \times \mathcal{%
M}_{2}\left( \mathbb{R}\right) \ $l'application $\varphi$ par : $\varphi \left( A,A'\right)
=\text{tr}\left( A^TA'\right)$, où $\text{tr}\left( A^TA'\right)$ désigne la trace du produit de la matrice $A^T$ par la matrice $A'$.\\ 
On admet que $\varphi $ est un produit scalaire sur $\mathcal{M}_{2}\left( 
\mathbb{R}\right)\ .$\\
\medskip
On note $\mathcal{F}=\left\{ \left( 
\begin{array}{cc}
a & b \\ 
-b & a%
\end{array}%
\right),\ \left( a,b\right) \in \mathbb{R}^{2}\right\}$.\\

\begin{enumerate}
\item Démontrer que $\mathcal{F}$ est un sous-espace vectoriel de $\mathcal{M}_{2}\left( \mathbb{R}\right) $.

\item D\'{e}terminer une base de $\mathcal{F}^{\perp }$.

\item D\'{e}terminer le projeté orthogonal de $J=\left( 
\begin{array}{cc}
1 & 1 \\ 
1 & 1%
\end{array}%
\right)$ sur $\mathcal{F}^{\perp }$ .
\item
Calculer la distance de $J$ à $\mathcal{F}$.
\end{enumerate}


\section*{EXERCICE 82 \: algèbre}

Soit $E$ un espace pr\'{e}hilbertien et $F$ un sous-espace vectoriel de $E$ de dimension finie $n>0$.\smallskip

On admet que, pour tout $x\in E$, il existe un élément unique $y_0$ de $F$ tel que $x-y_0$ soit orthogonal à $F$ et que la distance de $x$ à $F$ soit égale à  $\left\Vert x-y_0\right\Vert$.

Pour $A=\begin{pmatrix}
a & b \\ 
c & d%
\end{pmatrix}$ et $A'=\begin{pmatrix}
a^{\prime } & b^{\prime } \\ 
c^{\prime } & d^{\prime }%
\end{pmatrix}$,  on pose $\left(  A\ |\ A'\right) =aa^{\prime}+bb^{\prime}+cc^{\prime}+dd^{\prime}$.

\begin{enumerate}
\item Démontrer que  $\left( \,.\,|\,.\,\right) $ est un produit scalaire sur $\mathcal{M}_{2}\left( \mathbb{R}\right)$.

\item Calculer la distance de la matrice $A=\left( 
\begin{array}{cc}
1 & 0 \\ 
-1 & 2%
\end{array}%
\right)$ au sous-espace vectoriel $F$ des matrices triangulaires sup\'{e}rieures.
\end{enumerate}


\section*{Exercice 83 \: algèbre}



\addcontentsline{toc}{subsection}{\'Enoncé exercice 3}

Soit $u$ et $v$ deux endomorphismes d'un $\mathbb{R}$-espace vectoriel $E$.

\begin{enumerate}
\item Soit $\lambda$ un réel non nul. Prouver que si $\lambda$ est valeur propre de $u\circ v$, alors $\lambda$ est valeur propre de $v\circ u$.\\
\:\:\:\:
\item On considère, sur $E=\mathbb{R}\left[ X\right] $ les endomorphismes $u$ et $v$ définis par 
$u$:\:$P\longmapsto \displaystyle\int_{1}^{X}P $ et $v$:\:$P\longmapsto P'$ .\\
Déterminer $ \mathrm{Ker} (u\circ v)$ et $\mathrm{Ker} (v\circ u)$.
Le résultat de la question 1. reste-t-il vrai pour $\lambda=0$?\:\:\:\:
\item 
Si $E$ est de dimension finie, démontrer que le résultat de la première question reste vrai pour $\lambda=0$.\\
\textbf{Indication }: penser à utiliser le déterminant.\:\:\:\:

\end{enumerate}
\section*{EXERCICE 84 \: algèbre}



\begin{enumerate}
\item Donner la définition d'un argument d'un nombre complexe non nul (on ne demande ni l'interprétation géométrique, ni la démonstration de l'existence d'un tel nombre).
\item Soit $n\in{\mathbb{N}^{*}}$. Donner, en justifiant, les solutions dans  $\mathbb{C} $ de l'équation $z^{n}=1$ et préciser leur nombre.
\item En déduire, pour $n\in{\mathbb{N}^{*}}$, les solutions dans $\mathbb{C}$ de l'équation $\left( z+\mathrm{i}\right)^{n}=\left(z-\mathrm{i}\right)^{n}$ et démontrer que ce sont des nombres réels.
\end{enumerate}

\section*{EXERCICE 85\: algèbre}

\begin{enumerate}
\item Soient $n \in{\mathbb{N}^{*}}$, $P\in{\mathbb{R}_{n}\left[ X\right] }$ et $a \in{\mathbb{R}}$.
\begin{enumerate}
\item Donner sans démonstration, en utilisant la formule de Taylor, la décomposition de $P(X)$ dans la base 
$\left( 1, X-a, \left( X-a\right) ^{2}, \cdots, (X-a)^{n}\right) $.
\item 
Soit $r\in{\mathbb{N}^{*}}$. En déduire que :\\$a$ est une racine de $P$ d'ordre de multiplicité $r$  si et seulement si $P^{(r)}(a)\neq 0$
et $\forall k \in{\llbracket 0,r-1 \rrbracket}$  , $P^{(k)}(a)=0$.
\end{enumerate}
\item Déterminer deux réels $a$ et $b$ pour que $1$ soit racine double du polynôme $P=X^{5}+aX^{2}+bX$ et factoriser alors ce polynôme dans $\mathbb{R}\left[ X\right] $.
\end{enumerate} 


\section*{EXERCICE 86 \: algèbre}


\begin{enumerate}
\item

Soit $(a,b,p)\in{\mathbb{Z}^{3}}$.
Prouver que : si $p \wedge a = 1 $ et $p\wedge b=1$, alors $p\wedge (ab)=1$.
\item
Soit $p$ un nombre premier.
\begin{enumerate}

\item 
 Prouver que $\forall k \in{\llbracket 1,p-1\rrbracket}$, $p$ divise $\dbinom{p}{k}k!$\:\:  puis en déduire que $p$ divise $\dbinom{p}{k}$.



\item 
Prouver que:  $\forall n \in{\mathbb{N}}$,\: $n^{p}\equiv n \:\mod p $.\\
\textbf{Indication}: procéder par récurrence.
\item
En déduire, pour tout entier naturel $n$, que : $  p \:\text{ne divise pas }\: n \Longrightarrow\: n^{p-1}\equiv 1 \mod p$.

\end{enumerate}

\end{enumerate}


\section*{EXERCICE 87 \: algèbre }




Soient\ $a_{0},a_{1},\cdots ,a_{n}$ ,\: $n+1$ réels
deux à deux distincts.
\medskip

\begin{enumerate}
\item Montrer que si\ $b_{0},b_{1},\cdots ,b_{n}\;$sont$\;n+1\;$ réels \
quelconques, alors il existe un unique polyn\^{o}me $P\ $vérifiant
\begin{equation*}
\deg P\leqslant n \:\:\text{et}\:\:\forall i\in \llbracket 0,n\rrbracket,\:\:P\left(
a_{i}\right) =b_{i}.
\end{equation*}

\item Soit $k\in \llbracket 0,n \rrbracket $. \\
Expliciter ce polyn\^{o}me $P
$, que l'on notera $L_{k}$, lorsque :%
\begin{equation*}
\forall i\in \llbracket 0 ,n \rrbracket,  \ \ b_{i}=
\left\{ 
\begin{array}{lll}
0&\text{si}&i\neq k \\ 
1&\text{si}&i=k
\end{array}
\right. 
\end{equation*}

\item
 Prouver que $\forall p\in \llbracket 0 ,n \rrbracket$ ,\:\bigskip
 $\displaystyle\sum\limits_{k=0}^{n}a_{k}^{p}L_{k}=X^{p}$.

\end{enumerate}



\section*{EXERCICE 88 \: algèbre}

\begin{enumerate}
\item
Soit $E$ un $\mathbb{K}$-espace vectoriel ($\mathbb{K}=\mathbb{R}$ ou $\mathbb{C}$).\\
Soit $u\in\mathcal{L}(E)$.
Soit $P\in\mathbb{K}[X]$.\\
Prouver que si $P$ annule $u$ alors toute valeur propre de $u$ est racine de $P$.\\
\item
Soit $n\in\mathbb{N} $ tel que $n\geqslant 2$. On pose $E=\mathcal{M}_n\left( \mathbb{R}\right) $.\\
Soit $A=\left( a_{i,j}\right) _{\substack{ 1\leqslant i\leqslant n  \\ 1\leqslant j\leqslant n}}$ la matrice de $E$ définie par $a_{i,j}=\left\lbrace \begin{array}{l}
0\:\text{si}\:i=j\\
1\:\text{si}\:i\neq j\\
\end{array}\right. $ .\\

Soit $u\in\mathcal{L}\left(E\right) $ défini par : $\forall\:M\in E$, $u(M)=M+\mathrm{tr} (M)A$. \\
\begin{enumerate}
\item
Prouver que le polynôme $X^2-2X+1$ est annulateur de $u$.
\item
$u$ est-il diagonalisable?\\
Justifier votre réponse en utilisant deux méthodes (l'une avec, l'autre sans l'aide de la question 1.).

\end{enumerate}
\end{enumerate}
 

 \section*{EXERCICE 89 \: algèbre  }
 


                       
 Soit $n \in {\mathbb{N}}$ tel que $n\geqslant 2$. On pose $z = e^{\mathrm{i}\,\frac{2\pi}{n}}$.
\begin{enumerate}
 \item On suppose $k\in {\llbracket 1,n-1\rrbracket}$. \\Déterminer le module et un argument du complexe $\:z^k - 1$.
 \item On pose $\;S = \displaystyle\sum\limits_{k=0}^{n-1} \;\left|z^k - 1\right|$. Montrer que $\:S = \dfrac{2}{\tan\frac{\pi}{2n}}$.
 
 \end{enumerate}\vspace{0.2cm}
 



 \section*{EXERCICE 90 \: algèbre }


$\mathbb K$ désigne le corps des réels ou celui des complexes.\\
Soient $a_{1}, a_{2}, a_{3}$ trois scalaires distincts donnés de $\mathbb K$.

\begin{enumerate}
\item Montrer que $\Phi:\begin{array}[t]{ccc}
\mathbb K_{2}[X] & \longrightarrow & \mathbb K^{3} \hfill\\ 
P & \longmapsto & \big(P(a_{1}),P(a_{2}),P(a_{3})\big) \\
\end{array}$ est un isomorphisme d'espaces vectoriels.

\item On note $(e_{1}, e_{2}, e_{3})$ la base canonique de $\mathbb K^{3}$ et on pose  $\forall k\in\{1,2,3\}, \:L_{k}=\Phi^{-1}(e_{k})$.

\begin{enumerate}
\item Justifier que $(L_{1}, L_{2}, L_{3})$ est une base de $\mathbb K_{2}[X]$.

\item Exprimer les polynômes $L_{1}, L_{2}$ et $L_{3}$ en fonction de $a_{1}, a_{2}$ et $a_{3}$.
\end{enumerate}

\item Soit $P\in\mathbb K_{2}[X]$. Déterminer les coordonnées de $P$ dans la base $(L_{1}, L_{2}, L_{3})$.

\item  \textbf{Application} : on se place dans $\mathbb R^{2}$ muni d'un repère orthonormé et on considère les trois points $A(0,1), B(1,3), C(2,1)$. 

Déterminer une fonction polynomiale de degré 2 dont la courbe passe par les points $A$, $B$ et $C$.
\end{enumerate}


\section*{EXERCICE 91 \: algèbre}




 On considère la matrice $A=\begin{pmatrix}
0& 2 & -1 \\ 
-1& 3 &-1 \\ 
-1& 2 & 0
\end{pmatrix} \in\mathcal{M}_3(\mathbb{R}).$
\begin{enumerate}
\item Montrer  que $A$ n'admet qu'une seule valeur propre que l'on déterminera. 
\item La matrice $A$ est-elle inversible ? Est-elle diagonalisable ?
\item Déterminer, en  justifiant, le polynôme minimal de $A$. 
\item 
Soit $n\in\mathbb{N}$. Déterminer le reste de la division euclidienne de $X^n$ par $(X-1)^2$ et en déduire la valeur de $A^n$.
\end{enumerate}

\medskip

\section*{EXERCICE 92\: algèbre }

 
Soit $n\in\mathbb{N}^*$. On considère $E=\mathcal{M}_n(\mathbb{R})$ l'espace vectoriel des matrices carrées d'ordre $n$.\\
\medskip
On pose : $\forall(A,B)\in E^2$, $\left\langle A\:,B \right\rangle=\mathrm{tr}(A^TB)$ où tr désigne la trace et $A^T$ désigne la transposée de la matrice $A$.
\begin{enumerate}
\item
Prouver que  $\left\langle\:,\right\rangle$ est un produit scalaire sur $E$.
\item
On note $S_n(\mathbb{R})$ l'ensemble des matrices symétriques de $E$.\\
Une matrice $A$ de $E$ est dite antisymétrique lorsque $A^T=-A$.\\
On note $A_n(\mathbb{R})$ l'ensemble des matrices antisymétriques de $E$.\\
On admet que $S_n(\mathbb{R})$  et $A_n(\mathbb{R})$ sont des sous-espaces vectoriels de $E$.
\begin{enumerate}
\item
Prouver que $E=S_n(\mathbb{R})\oplus A_n(\mathbb{R})$.
\item
Prouver que $A_n(\mathbb{R})^\perp=S_n(\mathbb{R})$.
\end{enumerate}
\item
Soit $F$ l'ensemble des matrices diagonales de $E$.\\
Déterminer $F^\perp$.
\end{enumerate}

\medskip
\




\section*{EXERCICE 93 \: algèbre }



  Soit $E$ un espace vectoriel réel de dimension finie $n>0$ et $u\in{\mathcal{ L}(E)}$ tel que
   $u^3+u^2+u=0.$\\
   On notera $\mathrm{Id}$ l'application identité sur $E$.
   \begin{enumerate}
   \item Montrer que $\mathrm{Im }u \oplus \mathrm{Ker} u = E.$
   \item 
   \begin{enumerate}
   \item \'Enoncer le lemme des noyaux pour deux polynômes.
   \item En déduire que $\mathrm{Im }u =\mathrm{Ker} (u^2+u+\mathrm{Id})$.
   \end{enumerate}
  \item
  On suppose que $u$ est non bijectif.\\
Déterminer les valeurs propres de $u$. Justifier la réponse.
   
   \end{enumerate}
   \medskip
  \textbf{ Remarque}: les questions 1. , 2. et 3. peuvent être traitées  indépendamment les unes des autres.
    


\section*{EXERCICE 94 \: algèbre }



\begin{enumerate}
\item  En raisonnant par l'absurde, montrer que le système $(S)$: 
 $\left\{\begin{array}{lccl}
 x&\equiv&5& [6]\\
 x&\equiv&4&[8]
 \end{array}\right.$ n'a pas de solution $x$ appartenant à $\mathbb{Z}$
\item
  \begin{enumerate}
    \item \'Enoncer le théorème de Bézout dans $\mathbb{Z}$.\:\:\:\:
 \item
 Soit $a$ et $b$ deux entiers naturels premiers entre eux. \\
 Soit $c\in \mathbb{Z}$.\\
 \medskip
 Prouver que:  $\left( a|c\:\text{et} \:b|c\right)\Longleftrightarrow ab|c$. \:\:\:\
 \end{enumerate}
 \item 
 On considère le système $(S)$: 
 $\left\{\begin{array}{lccl}
           x&\equiv&6& [17]\\
            x&\equiv&5& [16]\\
 x&\equiv&4&[15]
 \end{array}\right.$ dans lequel l'inconnue $x$ appartient à $\mathbb{Z}$.\\
\begin{enumerate}
\item
Déterminer une solution particulière $x_0$ de $(S)$ dans $\mathbb{Z}$.
\item
\textit{Déduire des questions précédentes} la résolution dans  $\mathbb{Z}$ du système $(S)$.\\
On exprimera les solutions en fonction de la solution particulière $x_0$.\:\:\:\:
\end{enumerate}
\end{enumerate}

\pagebreak


\part*{BANQUE PROBABILIT\'ES}
\section*{EXERCICE 95 \:probabilités}


Une urne contient deux boules blanches et huit boules noires.\\
\begin{enumerate}
\item
Un joueur tire successivement, avec remise,  cinq boules dans cette urne.\\
Pour chaque boule blanche tirée, il gagne 2 points et pour chaque boule noire tirée, il perd 3 points.\\
On note $X$ la variable aléatoire représentant le nombre de boules blanches tirées.\\
On note $Y$ le nombre de points obtenus par le joueur sur une partie.
\begin{enumerate}
\item
Déterminer la loi de $X$, son espérance et sa variance.
\item
Déterminer la loi de $Y$, son espérance et sa variance.
\end{enumerate}
\item
Dans cette question, on suppose que l'on tire simultanément cinq boules dans l'urne.
\begin{enumerate}
\item 
Déterminer la loi de $X$.
\item
Déterminer la loi de $Y$.

\end{enumerate}
\end{enumerate}

\section*{EXERCICE 96 \:probabilités}



Soit $X$ une variable aléatoire à valeurs dans $\mathbb{N}$, de loi de probabilité donnée par : $\forall n\in \mathbb{N}$, $P(X=n)=p_n$.\\
  La fonction génératrice de $X$ est notée $G_X$ et elle est définie par $G_X(t)=E[t^X]=\displaystyle \sum_{n=0}^{+\infty}p_nt^n$.
  \begin{enumerate}
  \item
  Prouver que l'intervalle $\left] -1,1\right[ $ est inclus dans l'ensemble de définition de $G_X$.
\item
Soit $X_1$ et $X_2$ deux variables aléatoires indépendantes à valeurs dans $\mathbb{N}$.\\
On pose $S=X_1+X_2$.\\
Démontrer  que  $\forall t\in \left]-1,1 \right[ $, $G_S(t)=G_{X_1} (t)G_{X_2}(t)$:\\
\begin{enumerate}
\item
en utilisant le produit de Cauchy de deux séries entières.
\item
en utilisant uniquement la définition de la fonction génératrice.
\end{enumerate}

\textbf{Remarque}: on admettra, pour la question suivante,  que ce résultat est généralisable à $n$ variables aléatoires indépendantes à valeurs dans $\mathbb{N}$.
\item
Un sac contient quatre boules : une boule numérotée 0, deux boules numérotées 1 et une boule numérotée 2.\\
Soit $n\in\mathbb{N}^{*}$.
On effectue $n$ tirages  successifs, avec remise, d'une boule dans ce sac.\\
On note $S_n$ la somme des numéros tirés.\\
Soit $t\in \left]-1,1 \right[ $. \\
Déterminer $G_{S_n}(t)$ puis  en  déduire la loi de $S_n$.
  \end{enumerate}


\section*{EXERCICE 97 \:probabilités}


 Soit $(X,Y)$ un couple de variables aléatoires à valeurs dans $\mathbb{N}^2$ dont la loi est donnée par:\\
$\forall (j,k)\in {\mathbb{N}^2}$, 
$P\left( (X,Y)=(j,k)\right)=\dfrac{(j+k)\left( \dfrac{1}{2}\right) ^{j+k}}{\mathrm{e}\:j!\:k!}$.
\begin{enumerate}

\item
Déterminer les lois marginales de $X$ et de $Y$.\\
Les variables $X$ et $Y$ sont-elles indépendantes?
\item
Prouver que $E\left[ 2^{X+Y}\right] $ existe et la calculer.
\end{enumerate}

\section*{EXERCICE 98 \:probabilités}
Une secrétaire effectue, une première fois, un appel téléphonique vers $n$ correspondants distincts.\\
On admet que les $n$ appels constituent $n$ expériences indépendantes et que, pour chaque appel, la probabilité d'obtenir le correspondant demandé est $p$ ($p\in{\left]  0,1\right[ }$).\\
Soit $X$ la variable aléatoire représentant le nombre de correspondants obtenus.
\begin{enumerate}
\item Donner la loi de $X$. Justifier.
\item
La secrétaire rappelle une seconde fois, dans les mêmes conditions, chacun des $n-X$ correspondants qu'elle n'a pas pu joindre au cours de la première série d'appels.
On note $Y$ la variable aléatoire représentant le nombre de personnes jointes au cours de la seconde série d'appels.
\begin{enumerate}
\item
Soit $i\in \llbracket 0,n \rrbracket $.
Déterminer, pour  $k\in \mathbb{N}, $ $P(Y=k|X=i)$.
\item
Prouver que $Z=X+Y$ suit une loi binomiale dont on déterminera le paramètre.\\
\textbf{Indication} : on pourra utiliser, sans la prouver, l'égalité suivante: $\dbinom{n-i}{k-i}\dbinom{n}{i}=\dbinom{k}{i}\dbinom{n}{k}$.\\
\item
Déterminer l'espérance et la variance de $Z$.
\end{enumerate}
\end{enumerate}

\section*{EXERCICE 99 \:probabilités}

\begin{enumerate}
\item
Rappeler l'inégalité de Bienaymé-Tchebychev.
\item
 Soit $(Y_n)$ une suite de variables aléatoires  indépendantes, de même loi et et telle que $\forall n\in \mathbb{N}$, $Y_n\in L^2$.\\  On pose $S_n=\displaystyle\sum\limits_{k=1}^{n}Y_k$.\\

Prouver que: $\forall\:a\in \left] 0,+\infty\right[ $, $P\left( \left|\dfrac{S_n}{n}-E(Y_1)\right|\geqslant a\right) \leqslant\dfrac{V(Y_1)}{na^2}$.
\item \textbf{Application}\\
On effectue des tirages successifs, avec remise, d'une boule dans une urne contenant 2 boules rouges et 3 boules noires.\\
\`A partir de quel nombre de tirages peut-on garantir à plus de 95\% que la proportion de boules rouges obtenues restera comprise entre $0,35$ et $0,45$?\\
\textbf{Indication} : considérer la suite $(Y_i)$ de variables aléatoires de Bernoulli où $Y_i$ mesure l'issue du $i^{\text{ème}}$  tirage.

\end{enumerate}
 
 

\section*{EXERCICE 100 \:probabilités}
Soit $\lambda \in{\left] 0,+\infty\right[ }$.\\
Soit $X$ une variable aléatoire discrète à valeurs dans $\mathbb{N}^\ast$.\\
On suppose que $\forall n\in\mathbb{N}^\ast$, $P(X=n)=\dfrac{\lambda}{n(n+1)(n+2)} $.
\begin{enumerate}
\item Décomposer en éléments simples la fraction rationnelle $R$ définie par $R(x)=\dfrac{1}{x(x+1)(x+2)}$.
\item
Calculer $\lambda$.
\item
Prouver que $X$ admet une espérance, puis la calculer.
\item
$X$ admet-elle une variance? Justifier.

\end{enumerate}


\section*{EXERCICE 101 \:probabilités}


Dans une zone désertique, un animal  erre entre trois points d'eau $A$, $B$ et $C$.\\
\`A l'instant $t=0$, il se trouve au point A.\\
Quand il a épuisé l'eau du point où il se trouve, il part avec équiprobabilité rejoindre l'un des deux autres points d'eau.\\
 L'eau du point qu'il vient de quitter se régénère alors.\\
\medskip
Soit $n\in\mathbb{N}$.\\
On note $A_n$ l'événement «l'animal est en $A$  après son $n^{\text{ième}}$ trajet».\\
On note $B_n$ l'événement «l'animal est en $B$  après son $n^{\text{ième}}$ trajet».\\
On note $C_n$ l'événement «l'animal est en $C$  après son $n^{\text{ième}}$ trajet».\\
\medskip
On pose $P(A_n)=a_n$, $P(B_n)=b_n$ et $P(C_n)=c_n$.

\begin{enumerate}
\item 
\begin{enumerate}
\item
Exprimer, en le justifiant, $a_{n+1}$ en fonction de $a_n$, $b_n$ et $c_n$.
\item Exprimer, de même, $b_{n+1}$ et $c_{n+1}$ en fonction de $a_n$, $b_n$ et $c_n$.

\end{enumerate}

\item
\begin{spacing}{1,3}
On considère la  matrice $A=\left( {\begin{array}{ccc}
  0&\frac{1}{2}&\frac{1}{2} \\ 
   \frac{1}{2}&0& \frac{1}{2}\\
  \frac{1}{2}&\frac{1}{2}&0 \\ 
  
\end{array}} \right)$.\\
\end{spacing}
\begin{enumerate}


\item 
Justifier, sans calcul, que la matrice $A$ est diagonalisable.
\item
Prouver que $-\dfrac{1}{2}$ est valeur propre de $A$ et déterminer le sous-espace propre associé.
\item
Déterminer une matrice $P$  inversible  et une matrice $D$  diagonale de $\mathcal{M}_3(\mathbb{R})$ telles que $D=P^{-1}AP$.\\
\textbf{Remarque}: le calcul de $P^{-1}$ n'est pas demandé.
\end{enumerate}
\item
Montrer comment les résultats de la question 2. peuvent être utilisés pour calculer $a_n$, $b_n$ et $c_n$ en fonction de $n$.\\
\textbf{Remarque}: aucune expression finalisée de $a_n$, $b_n$ et $c_n$ n'est demandée.
\end{enumerate}
\bigskip


\section*{EXERCICE 102 \:probabilités}

Soit $N\in\mathbb{N}^*$.\\
Soit $p\in\left] 0,1\right[$. On pose $q=1-p$.\\
On considère $N$ variables aléatoires $X_1,X_2,\cdots,X_N$  définies sur un même espace probabilisé $\left(\Omega,\mathcal{A},P\right)$,  indépendantes et de même loi géométrique de paramètre $p$.
\begin{enumerate}
\item
Soit $i\in\llbracket1,N\rrbracket$. Soit $n\in{\mathbb{N}^*}$.\\
Déterminer $P(X_i\leqslant n)$, puis $P(X_i> n)$.
\item
On considère la variable aléatoire $Y$ définie par $Y=\underset{1\leqslant i\leqslant N}{\min}(X_i)$\\
c'est-à-dire $\forall \omega \in\Omega$, $Y(\omega)=\min\left(X_1(\omega),\cdots,X_N(\omega)\right)$, $\min$ désignant « le plus petit élément de ».
\begin{enumerate}
\item 
Soit $n\in{\mathbb{N}^*}$. Calculer $P(Y>n)$.\\ En déduire $P(Y\leqslant n)$, puis $P(Y=n)$.
\item
Reconnaître la loi de $Y$. En déduire $E(Y)$.
\end{enumerate}	

\end{enumerate}


\section*{EXERCICE 103 \:probabilités}




\textbf{Remarque}: les questions 1. et 2. sont indépendantes.\\
\medskip
Soit $(\Omega,\mathcal{A},P)$ un espace probabilisé.
\begin{enumerate}
\item 
\begin{enumerate}
\item
 Soit $\left( \lambda_1,\lambda_2 \right) \in \left( \left] 0,+\infty \right[\right) ^2$.\\
Soit $X_1$ et $X_2$ deux variables aléatoires définies sur $(\Omega,\mathcal{A},P)$.\\
On suppose que  $X_1$ et $X_2$ sont indépendantes et suivent des lois de Poisson, de paramètres respectifs  $\lambda_1$ et $\lambda_2$.\\
Déterminer la loi de $X_1+X_2$.\:\:\:\:
\item
En déduire l'espérance et la variance de $X_1+X_2$. \:\:\:\:
\end{enumerate}
\item 
Soit $p\in\left]  0,1\right]$. Soit $\lambda \in \left] 0,+\infty \right[$. \\ 
Soit $X$ et $Y$ deux variables aléatoires définies sur $(\Omega,\mathcal{A},P)$. \\
On suppose que  $Y$ suit une loi de Poisson de paramètre $\lambda$.\\
On suppose que $X(\Omega)=\mathbb{N}$ et que, pour tout $m\in\mathbb{N}$, la loi conditionnelle de $X$ sachant $(Y=m)$ est une loi binomiale de paramètre $(m,p)$.\\
Déterminer la loi de $X$.\:\:\:\:
\end{enumerate}


\section*{EXERCICE 104 \:probabilités}
Soit $n$ un entier naturel supérieur ou égal à 3.\\
On dispose de $n$ boules numérotées de $1$ à $n$ et d'une boîte formée de trois compartiments identiques également numérotés de 1 à 3.\\
On lance simultanément les $n$ boules.\\
Elles viennent toutes se ranger aléatoirement dans les  3 compartiments.\\
Chaque compartiment peut éventuellement contenir les $n$ boules.\\
On note $X$ la variable aléatoire qui à chaque expérience aléatoire fait correspondre le nombre de compartiments restés vides.%pourquoi indexer X par n le nombre de boules?%
\begin{enumerate}
\item
Préciser les valeurs prises par $X$.
\item
\begin{enumerate}
\item
Déterminer la probabilité $P(X=2)$.
\item
Finir de déterminer la loi de probabilité de $X$.
\end{enumerate}
\item
\begin{enumerate}
\item
Calculer $E(X)$.
\item
Déterminer $\lim\limits_{n\to +\infty}^{}E(X)$. Interpréter ce résultat.
\end{enumerate}
\end{enumerate}

\section*{EXERCICE 105 \:probabilités}

\begin{enumerate}
\item
\'Enoncer et démontrer la formule de Bayes pour un système complet d'événements.
\item
On dispose de 100 dés dont 25 sont pipés (c'est-à-dire truqués).\\
Pour chaque dé pipé, la probabilité d'obtenir le chiffre 6 lors d'un lancer vaut $\dfrac{1}{2}$.
\begin{enumerate}


\item
On tire un dé au hasard parmi les 100 dés. On lance ce dé  et on obtient le chiffre 6.\\
Quelle est la probabilité que ce dé soit pipé?
\item 
Soit $n\in\mathbb{N}^*.$ \\
On tire un dé au hasard parmi les 100 dés. On lance ce dé $n$ fois et on obtient $n$ fois le chiffre $6$.\\
Quelle est la probabilité $p_n$ que ce dé soit pipé?
\item
Déterminer $\lim\limits_{n\to+\infty}^{}p_n$. Interpréter ce résultat.
\end{enumerate}
\end{enumerate}

\section*{EXERCICE 106 \:probabilités}


 $X$ et $Y$ sont deux variables aléatoires indépendantes et à valeurs dans $\mathbb{N}$.\\
Elles suivent la même loi définie par:  
$\forall\:k\in\mathbb{N}$, $P(X=k)=P(Y=k)=pq^k$ où    $p \in \left] 0,1\right[ $ et $q=1-p$.\\
On considère alors les variables $U$ et $V$ définies par $U=\sup(X,Y)$ et $V=\inf(X,Y)$.
\begin{enumerate}
\item
Déterminer la loi du couple $(U,V)$.
\item
Déterminer la loi marginale de $U$.\\
On admet que $V(\Omega)=\mathbb{N}$ et que, $\forall\:n\in\mathbb{N}$, $P(V=n)=pq^{2n}(1+q)$.
\item
Prouver que $W=V+1$ suit une loi géométrique.\\
En déduire l'espérance de $V$.
\item
$U$ et $V$ sont-elles indépendantes?

\end{enumerate}

\section*{EXERCICE 107 \:probabilités}

On dispose de deux urnes $U_1$ et $U_2$.\\
L'urne $U_1$ contient deux boules blanches et trois boules noires.\\
L'urne $U_2$ contient quatre boules blanches et trois boules noires.\\
On effectue des tirages successifs dans les conditions suivantes:\\
on choisit une urne au hasard et on tire une boule dans l'urne choisie.\\
On note sa couleur et on la remet dans l'urne d'où elle provient.\\
Si la boule tirée était blanche, le tirage suivant se fait dans l'urne $U_1$.\\
Sinon le tirage suivant se fait dans l'urne $U_2$.\\
Pour tout  $n\in\mathbb{N}^*$, on note $B_n$ l'événement « la boule tirée au $n^{\text{ième}}$ tirage est blanche »
et on pose $p_n=P(B_n)$.
\begin{enumerate}
\item
Calculer $p_1$.
\item
Prouver que : $\forall\:n\in\mathbb{N}^*$, $p_{n+1}=-\dfrac{6}{35}p_n+\dfrac{4}{7}$.
\item
En déduire, pour tout entier naturel $n$ non nul, la valeur de $p_n$.
\end{enumerate}

\section*{EXERCICE 108 \:probabilités}


Soient $X$ et $Y$ deux variables aléatoires définies sur un même espace probabilisé $(\Omega,\mathcal{A},P)$ et à valeurs dans $\mathbb{N}$.\\
On suppose que la loi du couple $(X,Y)$ est donnée par:\\
 $\forall\:(i,j)\in\mathbb{N}^2$, $P((X=i)\cap (Y=j))=\dfrac{1}{\mathrm{e}\:2^{i+1}j!}$
\begin{enumerate}


\item
Déterminer les lois de $X$ et de $Y$.
\item
\begin{enumerate}
\item 
Prouver que $1+X$ suit une loi géométrique et en déduire l'espérance et la variance de $X$.
\item
Déterminer l'espérance et la variance de $Y$.
\end{enumerate}

\item
Les variables $X$ et $Y$ sont-elles indépendantes?
\item
Calculer $P(X=Y)$.
\end{enumerate}

\section*{EXERCICE 109 \:probabilités}

Soit $n\in\mathbb{N}^*$.
Une urne contient $n$ boules blanches numérotées de 1 à $n$ et deux boules noires numérotées 1 et 2.\\
On effectue le tirage une à une, sans remise, de toutes les boules de l'urne.\\
On suppose que tous les tirages sont équiprobables.\\
On note $X$ la variable aléatoire égale au rang d'apparition de la première boule blanche.\\
On note $Y$ la variable aléatoire égale au rang d'apparition de la première boule numérotée 1.\\
\begin{enumerate}
\item 
Déterminer la loi de $X$.
\item
Déterminer la loi de $Y$.



\end{enumerate}  

\section*{EXERCICE 110 \:probabilités}

Soit $(\Omega,\mathcal{A},P)$ un espace probabilisé.\\

\begin{enumerate}
\item 
Soit $X$ une variable aléatoire définie sur  $(\Omega,\mathcal{A},P)$ et à valeurs dans $\mathbb{N}$.\\
On considère la série entière  $\displaystyle\sum t^nP(X=n)$ de variable réelle $t$.\\
On note $R_X$ son rayon de convergence.\\
\begin{enumerate}
\item 
Prouver que $R_X\geqslant 1$.\:\:\:\:\\
On pose  $G_X(t)=\displaystyle\sum\limits_ {n=0}^{+\infty}t^nP(X=n)$ et on  note $D_{G_X}$ l'ensemble de définition de $G_X$.\\
Justifier que   $\left[-1,1 \right] \subset D_{G_X}$.\:\:\:\:\\
 
Pour tout réel $t$ fixé de $\left[-1,1 \right]$, exprimer $G_X(t)$ sous forme d'une espérance.\:\:\:\:
\item
Soit $k\in\mathbb{N}$.
Exprimer, en justifiant la réponse, $P(X=k)$ en fonction de  $G_X^{(k)}(0)$.\:\:\:\:

\end{enumerate}
\item 
\begin{enumerate}
\item
 On suppose que $X$ suit une loi de Poisson de paramètre $\lambda$.\\
 Déterminer $D_{G_X}$ et, pour tout $\:t\in D_{G_X}$, calculer $G_X(t)$.\:\:\:\:
 

 \item
 Soit $X$ et $Y$ deux variables aléatoires définies sur un même espace probabilisé, indépendantes  et suivant des lois de Poisson de paramètres respectifs $\lambda_1$ et $\lambda_2$.\\
 Déterminer, en utilisant les questions précédentes, la loi de $X+Y$.\:\:\:\:
 \end{enumerate}
\end{enumerate}

\section*{EXERCICE 111 \:probabilités}

On admet, dans cet exercice, que:
 $\forall\:q\in \mathbb{N}$, $\forall \:x\in \left] -1,1\right[$, $\displaystyle\sum\limits_{k\geqslant q}^{}\dbinom {k}{q}x^{k-q}$ converge et  que  $\displaystyle\sum\limits_{k=q}^{+\infty}\dbinom{k}{q}x^{k-q}=\dfrac{1}{(1-x)^{q+1}}.$\\
 \bigskip
Soit $p\in \left] 0,1\right[ $.\\
Soit $(\Omega,\mathcal{A},P)$ un espace probabilisé.\\
Soit $X$ et $Y$ deux variables aléatoires définies sur  $(\Omega,\mathcal{A},P)$ et à valeurs dans $\mathbb{N}$.\\
On suppose que la loi de probabilité du  couple $(X,Y)$ est donnée par:\\
\medskip
$\forall\:(k,n)\in\mathbb{N}^2$, $P((X=k)\cap (Y=n))=\left\lbrace \begin{array}{l}
\dbinom{n}{k}\left( \dfrac{1}{2}\right)^np(1-p)^n\:\text{si}\:k\leqslant n\\
0\:\text{sinon} 
\end{array}
\right.$
\begin{enumerate}
\item
Vérifier qu'il s'agit bien d'une loi de probabilité.
\item
\begin{enumerate}
\item
Déterminer la loi de $Y$.
\item
Prouver que $1+Y$ suit une loi géométrique.
\item
Déterminer l'espérance de $Y$.
\end{enumerate}
\item
Déterminer la loi de $X$.


\end{enumerate}

\section*{EXERCICE 112 \:probabilités}


Soit $n\in\mathbb{N}^*$ et $E$ un ensemble possédant $n$ éléments.\\
On désigne par $\mathcal{P}(E)$ l'ensemble des parties de $E$.
\medskip
\begin{enumerate}
\item
Déterminer le nombre $a$ de couples $(A,B)\in \left(\mathcal{P}(E) \right)^2$ tels que $A\subset B $.
\item
Déterminer le nombre $b$ de couples $(A,B)\in \left(\mathcal{P}(E) \right)^2$ tels que $A\cap B=\emptyset $.
\item
Déterminer le nombre $c$ de triplets  $(A,B,C)\in \left(\mathcal{P}(E) \right)^3$ tels que $A$, $B$ et $C$ soient deux à deux disjoints et vérifient $A\cup B\cup C=E$.

 
\end{enumerate}

   

 





\end{flushleft}
\end{document}
